\documentclass[11pt]{article}
\usepackage[T1]{fontenc}
\usepackage{lmodern}
\usepackage[margin=1.05in]{geometry}
\usepackage{amsmath,amssymb,amsthm,mathtools}
\usepackage{microtype}
\usepackage{enumitem}
\usepackage{xcolor}
\PassOptionsToPackage{hyphens}{url}
\usepackage{hyperref}
\hypersetup{colorlinks=true,linkcolor=blue!45!black,citecolor=blue!45!black,urlcolor=blue!45!black,pdftitle={A smooth three-torus diffeomorphism with simple Lebesgue spectrum},pdfauthor={OpenAI}}
\newcommand{\T}{\mathbb T}
\newcommand{\R}{\mathbb R}
\newcommand{\Z}{\mathbb Z}
\newcommand{\N}{\mathbb N}
\newcommand{\C}{\mathbb C}
\newcommand{\eps}{\varepsilon}
\newcommand{\1}{\mathbf 1}
\newcommand{\dd}{\,\mathrm d}
\DeclareMathOperator{\ee}{e}
\DeclareMathOperator{\Rea}{Re}
\DeclareMathOperator{\Ima}{Im}
\DeclareMathOperator{\supp}{supp}
\DeclareMathOperator{\dist}{dist}
\DeclareMathOperator{\Leb}{Leb}
\newcommand{\calH}{\mathcal H}
\newcommand{\calC}{\mathcal C}
\newtheorem{theorem}{Theorem}[section]
\newtheorem{lemma}[theorem]{Lemma}
\newtheorem{proposition}[theorem]{Proposition}
\newtheorem{corollary}[theorem]{Corollary}
\theoremstyle{definition}
\newtheorem{definition}[theorem]{Definition}
\theoremstyle{remark}
\newtheorem{remark}[theorem]{Remark}
\numberwithin{equation}{section}
\setlist{itemsep=2pt,topsep=5pt}
\setlength{\emergencystretch}{1em}
\title{A smooth three-torus diffeomorphism with simple Lebesgue spectrum}
\author{OpenAI}
\date{September 23, 2026}
\begin{document}
\maketitle
\begin{abstract}
We solve the probability-preserving form of Banach's simple Lebesgue-spectrum problem in smooth dynamics. We construct a $C^\infty$ diffeomorphism of the three-torus that preserves standard volume and whose Koopman operator has simple Lebesgue spectrum on the entire mean-zero $L^2$ space.
\end{abstract}
\tableofcontents
\section{Introduction}
\label{sec:introduction}
Let $\T=\R/\Z$, let $\mu$ be normalized Lebesgue measure on $\T^3$,
and put $\ee(t)=\exp(2\pi i t)$. For an invertible measure-preserving
transformation $T$, we use the convention $U_Tg=g\circ T$.
The closed subspace of complex square-integrable functions of integral zero
is denoted by $L^2_0(\T^3,\mu)$.

On a probability space the constants form an invariant subspace, so the
probability-preserving form of Banach's simple Lebesgue-spectrum problem
asks for an invertible transformation whose Koopman operator on the
\emph{entire} orthogonal complement of the constants is unitarily
equivalent to multiplication by $w$ on $L^2(S^1,m)$, where $m$ is
normalized Haar measure. Equivalently, that restriction has a bilateral
orthonormal orbit which is complete. We realize this spectral model with
the standard smooth volume on a compact manifold.

\begin{theorem}\label{thm:main}
There exist a $C^\infty$ volume-preserving diffeomorphism
$T\colon\T^3\to\T^3$ and a real function $f\in L^2_0(\T^3,\mu)$ such that
\[
                  \{f\circ T^n:n\in\Z\}
\]
is an orthonormal basis of $L^2_0(\T^3,\mu)$.
Consequently $T$ is ergodic, and $U_T$ on this space is unitarily equivalent
to multiplication by the coordinate $w$ on $L^2(S^1,m)$.
\end{theorem}

The invariant measure in Theorem~\ref{thm:main} is the standard smooth
volume itself. The generator $f$ is only required to belong to $L^2$.
In particular, the assertion concerns the whole mean-zero space, including
both its spectral type and its multiplicity.
Section~\ref{sec:consequences} derives mixing, zero entropy, and vanishing
Lyapunov exponents. Mixing of every order uses, in addition, the separate
multiple-mixing theorem of \cite{MultipleMixing2026}.

The formulation above is distinct from the historical real-line question
recorded by Ulam and attributed to Banach \cite[\S 6, p.~76]{Ulam1960}.
Rokhlin asked for ergodic automorphisms with simple, or at least
finite-multiplicity, Lebesgue spectrum
\cite[\S 4, no.~7, p.~107]{Rokhlin1949}. Helson and Parry constructed
cocycles over aperiodic probability-preserving transformations whose
associated unitary operators have Lebesgue spectrum
\cite[p.~196, Theorem~(a)]{HelsonParry1978}. Mathew and Nadkarni constructed a
probability-preserving transformation with a Lebesgue \emph{component}
of multiplicity two \cite{MathewNadkarni1984}. Guenais connected the
Morse-cocycle question with flat polynomials and constructed a countable
abelian group action with simple spectrum of mixed discrete and Haar
type \cite{Guenais1999}. These results distinguish a Lebesgue component,
or a related group action, from one bilateral Lebesgue orbit spanning
the whole centered space.

Simple spectrum in smooth dynamics was already achieved by Fayad, who
constructed $C^\infty$ volume-preserving torus diffeomorphisms with
simple spectrum \cite{Fayad2001}; that result did not assert Lebesgue
spectral type.
In another setting, an infinite-measure conservative transformation
with simple Lebesgue spectrum is presented in
\cite[Theorem~2.2]{ElAbdalaoui2023}; the probability-space conclusion
does not follow from it. Theorem~\ref{thm:main} gives that conclusion
for the full mean-zero space of a smooth probability-preserving map.

Related existence results in continuous time have different spectral
consequences. Prikhod'ko constructs a finite-measure flow with simple
Lebesgue spectrum \cite[Theorem~3]{Prikhodko2020}. In the smooth category,
Fayad, Forni and Kanigowski obtain a real analytic conservative flow on
$\T^2$ with one singularity and countably infinite Lebesgue multiplicity
\cite[Theorem~1]{FFK2021}; Abdedou, Fayad and Kessi obtain real analytic
reparametrizations of minimal translation flows on $\T^5$ with the same
multiplicity \cite[Theorem~1.1]{AFK2023}. Taking a nonzero time map of a
simple Lebesgue flow yields countably infinite multiplicity: its spectral
model is multiplication by $\ee(t\xi)$ on $L^2(\R,\dd\xi)$, and the
substitution $\xi=(x+k)/|t|$, with $x\in[0,1)$ and $k\in\Z$, gives one
Lebesgue copy for each $k$. Thus these continuous-time results do not
supply Theorem~\ref{thm:main}.

Successive smooth changes of coordinates have a substantial history in
smooth dynamics, including the approximation-by-conjugacy constructions
of Anosov and Katok \cite{AnosovKatok1970}. We use explicit
volume-preserving passages between integrable twists rather than an
external conjugacy theorem. Their quantitative spectral-density
transport, and the finite Fourier signals built on it, are proved below.

\subsection*{The construction}
The approximating maps are integrable twists. In a volume-preserving
coordinate system on the torus, they have the form
\[
 S(\theta,u,a)=(\theta+\phi(u),u,a),
\]
where $\phi$ is a smooth real circle function that is strictly increasing
on an arc whose complement can be made arbitrarily short. A finite angular
Fourier expansion supported on that arc has an explicit absolutely
continuous spectral measure. We refer to its individual harmonics as
\emph{packets}, and keep a finite label, or \emph{color}, on disjoint groups
of packets. The colors carry the coefficients of an observable that we
want to approximate by translates of a single vector.

There are three steps. A slow passage in the coordinate $u$ replaces the
twist by an arbitrarily close twist whose angular coordinate is the old
$u$. Stationary phase locates the new angular indices and determines the
spectral density at a much finer scale. The spare coordinate $a$ supplies
a small adjustable variation in the speed (the \emph{jitter}). A large integer added to the
lift of $\phi$ separates the angular indices of different packets without
changing the old map. This separation is what allows later perturbations
to distinguish colors even when their temporal spectra overlap.

Next, a positive Fourier construction and a finite prime sieve produce a
small sinusoidal bias in each color density. Its coefficient depends on
the normalized angular index. A neutral passage between successive biased
passages makes residue classes carry equal shares of each color. The
bias strength can be chosen before the angular scale grows; the passage
parameter is then free to enforce any required smooth or spectral
accuracy.

Finally, many small biases are grouped into blocks. If color $r$ has
spectral density $p_r$ and target coefficient $b_r$, the center at a
temporal frequency $z$ is the density-weighted mean
\[
 p(z)=\sum_r p_r(z),\qquad
 \beta(z)=\frac{\sum_r b_rp_r(z)}{p(z)}\quad(p(z)>0).
\]
Centering each block at a smooth approximation to this current mean
makes the loss in $\int\sqrt p$ quadratic
in the block's strength budget, while the same block provides an estimator
of the color label. Repetition makes prediction arbitrarily accurate with
arbitrarily little loss in $\int\sqrt p$. This separation of the two
estimates is the central finite construction. When the vector's norm is
close to one, keeping $\int\sqrt p$ close to one makes whitening a small
change in physical $L^2$, while restoring the spectral density exactly
to one. A diagonal limit preserves the orthonormal orbit and makes its
span contain every member of a dense family of observables.

\subsection*{Order of the proof}
Section~\ref{sec:passages} proves the analytic passage rules, including the
comparison of jitter laws along a rapidly oscillating phase.
Section~\ref{sec:signals} constructs the positive laws that transmit the
small biases. Section~\ref{sec:blocks} proves the finite prediction result,
and Section~\ref{sec:completion} prepares arbitrary observables and takes
the smooth limit. Section~\ref{sec:consequences} derives the dynamical
consequences. Every passage acts on finite data. Its integer parameter
is chosen last, after the current coordinate chart, cutoffs, jitter,
and all desired error bounds have been fixed. No uniform bound on the
number of packets or on the derivatives of successive coordinate charts
is assumed.

\section{Smooth passages and spectral densities}\label{sec:passages}

A volume-preserving coordinate system, also called a chart below, means
a smooth volume-preserving identification of the physical torus with $\T^3$.
All coordinate systems in the construction preserve normalized volume.
An intermediate transformation has the form
\begin{equation}\label{eq:twist}
 S(\theta,u,a)=(\theta+\phi(u),u,a),
 \qquad (\theta,u,a)\in\T^3,
\end{equation}
where $\phi\in C^\infty(\T;\R)$ has $\phi'>0$ on an open arc
$I$. We call $I$ a good arc. Derivatives on an arc are understood in an
increasing real lift of that arc.

A packet is a function $A(u,a)\ee(m\theta)$, where
$m\in\Z\setminus\{0\}$ and $A$ is smooth with compact support in
$I\times\T$. A finite packet list is disjoint if, for any two entries with
the same $m$, their coefficient supports are disjoint. Entries with
different $m$ need no support restriction. Packets may be grouped into
colors; we write $F_r$ for the sum of color $r$. Disjointness implies
\begin{equation}\label{eq:packet-orthogonality}
 \langle U_S^n F_r,U_S^k F_l\rangle=0
 \quad(r\ne l,\ n,k\in\Z).
\end{equation}
Indeed $U_S^n(A\ee(m\theta))=\ee(nm\phi(u))A\ee(m\theta)$,
and either angular orthogonality or disjoint coefficient supports applies.
The spectral measure of one packet is
\begin{equation}\label{eq:packet-spectral-measure}
 (m\phi\bmod1)_*\bigl(|A(u,a)|^2\dd u\dd a\bigr).
\end{equation}
It has a smooth density because its coefficient support is compactly
inside the good arc. We use the inner product linear in its first
variable, so a vector of density $p$ has
$\langle U_S^nF,F\rangle=\int_\T\ee(nz)p(z)\dd z$.
A measurable spectral multiplier $v(z)$ acts on its
coefficient by multiplication by $v(m\phi(u)\bmod1)$, whenever the result
belongs to $L^2$. These statements follow directly from the displayed
formula for $U_S^n$, first for trigonometric polynomials and then by
$L^2$ approximation. The pushforward formula, absolute continuity,
and multiplier rule also apply to nonsmooth coefficients in $L^2$
supported in the good arc; their densities need not be smooth.

For later comparisons, if two vectors have spectral densities $p_F,p_G$
for the same transformation, then
\begin{equation}\label{eq:density-norm-bound}
 \|p_F-p_G\|_1
 \le (\|F\|_2+\|G\|_2)\|F-G\|_2.
\end{equation}
To see this, test their spectral measures against an arbitrary function
of absolute value at most one and use the contraction property of its
spectral multiplier. The total variation norm here and below is the
$L^1$ norm of densities, without a factor $1/2$.

\begin{lemma}[Passage]\label{lem:passage}
Fix \eqref{eq:twist}, its physical volume-preserving coordinate chart,
and a finite disjoint packet list. Choose a positive integer $c$ so large
that the compact ranges $m(c+\phi(\T))$, for the distinct old indices $m$,
are pairwise disjoint and avoid zero. Set $V=c+\phi$.
Fix $J\in C^\infty(\T;\R)$ with $J'>0$ on an open arc $I'$ and fix
$0\le\psi\le1$ in $C_c^\infty(I')$.
For integers $N\to\infty$, choose $\Delta_N$ so that
\begin{equation}\label{eq:passage-rounding}
 |\Delta_N|\le N^{-1},\qquad
 k_N:=N\left(\int_\T V(u)\dd u+\Delta_N\right)\in\Z,
 \qquad H_N'=N(V+\Delta_N).
\end{equation}
Put $\sigma_N(a)=N^{-1}+N^{-2}J(a)$. Then
\begin{equation}\label{eq:passage-map}
 T_N(\theta,u,a)=
 \bigl(\theta+H_N(u+\sigma_N(a))-H_N(u),\,
 u+\sigma_N(a),\,a\bigr)
\end{equation}
is a volume-preserving smooth torus diffeomorphism, and both $T_N$ and
$T_N^{-1}$ converge smoothly to $S$ and $S^{-1}$, also in the fixed
physical coordinates. In the new coordinates
\begin{equation}\label{eq:passage-chart}
 (\theta',u',a')=C_N(\theta,u,a)
             =(u,a,\theta-H_N(u)),
\end{equation}
it is the twist of speed $\sigma_N$ and good arc $I'$.

Here are its packet and density properties. For an old packet indexed
by $\alpha$, with index $m_\alpha$ and amplitude $A_\alpha$, define
\begin{equation}\label{eq:stationary-image}
 E_\alpha=\{(m_\alpha V(u),a):(u,a)\in\supp A_\alpha\}.
\end{equation}
One may choose pairwise disjoint, arbitrarily small fixed neighborhoods
of these compact sets, and smooth cutoffs $\chi_\alpha(s,a)$ supported
in those neighborhoods and equal to one near $E_\alpha$. Retain the
finitely many new packets
\begin{align}
 G_{\alpha,j}(\theta',u',a')
   &=\ee(j\theta')\ee(m_\alpha a')
       \psi(u')\chi_\alpha(j/N,u')C_{\alpha,j}(u'),
       \label{eq:retained-packets}\\
 C_{\alpha,j}(a)
   &=\int_\T A_\alpha(u,a)\ee(m_\alpha H_N(u)-ju)\dd u.
       \label{eq:exact-passage-coefficient}
\end{align}
All their angular indices are nonzero for sufficiently large $N$.
For each $\alpha$,
\begin{equation}\label{eq:transport-norm}
 \left\|\sum_jG_{\alpha,j}\circ C_N
       -\psi(a)A_\alpha(u,a)\ee(m_\alpha\theta)\right\|_2=o(1).
\end{equation}
They are disjoint as packets, and their moduli are independent of the
new spare coordinate $a'$. Every constant attached to an old packet is
preserved by this linear transport and these cutoffs.

For a single old packet write, with zero extension outside its branch,
\begin{equation}\label{eq:W-definition}
 W(s,a)=\frac{|A(u,a)|^2}{|m\phi'(u)|},\qquad s=mV(u).
\end{equation}
Let $K(s,y)$ be the density of the pushforward of
$W(s,a)\psi(a)^2\dd a$ under $a\mapsto sJ(a)$. It is continuous and
compactly supported in $(s,y)$, with its $s$ support away from zero.
For any fixed positive integer $P$ and residue $h\pmod P$, the unfolded
temporal density of the retained packets with $j\equiv h\pmod P$ differs
in $L^1(\R)$ by $o(1)$ from
\begin{equation}\label{eq:fine-density-unfolded}
 Q_{N,h}(t)=\sum_{k\in\Z}K(t,Nt-h-kP).
\end{equation}
Consequently their density on $\T$ differs in $L^1(\T)$ by $o(1)$ from
$\sum_{l\in\Z}Q_{N,h}(z+l)$. The assertions hold after summing over
any of the finitely many colors. All $o(1)$ terms refer to the old data,
$c,J,\psi,P$, and support neighborhoods fixed before $N$ tends to
infinity. In particular, $N$ may be required to exceed any prescribed
bound.
\end{lemma}

\begin{proof}
\emph{Geometry and smooth convergence.}
Choose $k_N$ by rounding $N\int V$ to an integer and integrate
\eqref{eq:passage-rounding}. Then $H_N(u+1)-H_N(u)=k_N$, so the map
$u\mapsto H_N(u)\bmod1$ is globally defined. The map $C_N$ is a smooth
torus diffeomorphism of Jacobian one, with inverse
$(\theta',u',a')\mapsto(a'+H_N(\theta'),\theta',u')$.
Direct substitution gives
\[
 C_NT_NC_N^{-1}(\theta',u',a')
       =(\theta'+\sigma_N(u'),u',a').
\]
This proves the global assertions, except convergence. The integral
identity
\[
 H_N(u+\sigma_N(a))-H_N(u)
       =N\int_0^{\sigma_N(a)}(V(u+t)+\Delta_N)\dd t
       =V(u)+O_{C^k}(N^{-1})
\]
holds for each fixed $k$, with constants depending on the fixed $c,J$
and $\phi$. Since $c$ is an integer, addition of $V$ and of $\phi$ agree
on $\T$. The inverse of \eqref{eq:passage-map} is
\[
 (\theta,u,a)\longmapsto
 (\theta-H_N(u)+H_N(u-\sigma_N(a)),u-\sigma_N(a),a),
\]
and the same estimate proves inverse convergence. Composing with a
fixed physical chart and its inverse preserves convergence in every
fixed smooth norm. No bound uniform over future coordinate charts is
being asserted.

\emph{Fourier coefficients.}
Work on a compact
subarc of $I$ containing all old coefficient supports in its interior.
For $s=j/N$ the phase in \eqref{eq:exact-passage-coefficient}, divided
by $N$, has critical equation
\begin{equation}\label{eq:critical-equation}
 m(V(u)+\Delta_N)=s
\end{equation}
and second derivative $m\phi'(u)$. It has at most one critical point
on this arc. With $u_*=u_*(s-m\Delta_N)$ on that branch, set
\[
 L_{m,j}(a)=
 \frac{A(u_*,a)}{\sqrt{N|m\phi'(u_*)|}}
 \ee\left(mH_N(u_*)-ju_*+\frac{\operatorname{sgn}(m)}8\right),
\]
and set this to zero wherever the amplitude at the critical point is
zero or there is no point on the branch. Uniformly for $j/N$ in any
fixed compact interval and $a\in\T$,
\begin{equation}\label{eq:stationary-uniform}
 C_{m,j}(a)-L_{m,j}(a)=o(N^{-1/2}).
\end{equation}

For general stationary-phase background see \cite{Hormander2003};
we give the compact-support one-dimensional argument needed here.
Enlarge the support arc twice inside $I$. Critical points
near the original support then range over a fixed compact interval on
which the absolute second derivative is bounded below. A smooth change
of variable, depending smoothly on the critical point, reduces the
divided phase to its critical value plus
$\operatorname{sgn}(m)y^2/2$; its Jacobian at zero is
$|m\phi'(u_*)|^{-1/2}$. After putting $v=\sqrt N y$, the integral on
$|v|\le D$ converges uniformly to the constant amplitude times the
quadratic exponential integral there. The remaining local tails are
$O(N^{-1/2}/D)$, uniformly in the parameters, by integrating against
$\ee(\operatorname{sgn}(m)v^2/2)$: the rescaled amplitudes have uniformly
bounded suprema and total variation. The part away from the critical
point is $O(N^{-1})$ by integration by parts. Finally
\[
 \lim_{D\to\infty}\int_{-D}^D\ee(\pm v^2/2)\dd v
          =\exp(\pm i\pi/4);
\]
this follows by inserting a Gaussian factor, evaluating the Gaussian
integral, and removing the factor using the same tail bound.
Taking first $N\to\infty$, then $D\to\infty$, proves
\eqref{eq:stationary-uniform}. When the potential critical point is
outside the enlarged support interval, the integral is uniformly
nonstationary. This also covers zero amplitudes and the ends of the
branch.

There are $O(N)$ indices in a fixed bounded $j/N$ range. Outside a
sufficiently large such range, repeated integration by parts gives,
for any fixed positive integer $q$,
$|C_{m,j}(a)|\le C_q|j|^{-q}$ uniformly in $a$ and $N$.
Indeed the phase derivative has size comparable to $|j|$, and its
higher derivatives are $O(N)$ with $N=O(|j|)$ there. Thus
\eqref{eq:stationary-uniform} implies
\[
 \sum_j\int_\T |C_{m,j}-L_{m,j}|^2\dd a=o(1).
\]
The $L^2$ norm of the leading array is bounded by a Riemann sum.
Cauchy--Schwarz therefore gives the useful mass estimate
\begin{equation}\label{eq:summed-stationary-mass}
 \sum_j\int_\T
 \left||C_{m,j}(a)|^2
       -\frac1N W(j/N-m\Delta_N,a)\right|\dd a=o(1).
\end{equation}

\emph{Cutoffs and transport.}
The map $(u,a)\mapsto(mV(u),a)$ is injective on $I\times\T$.
It takes disjoint compact supports with the same $m$ to disjoint
compact sets. Our choice of $c$ separates the images for different
$m$. The stated cutoffs $\chi_\alpha$ therefore exist. For large $N$
they equal one on the support of the corresponding leading array,
since its $s$ coordinate is shifted by only $m_\alpha\Delta_N$.
Because each cutoff is one on the shifted leading-array support, the
deleted exact coefficient mass is bounded by the summed squared
$C_{m,j}-L_{m,j}$ error, which tends to zero by the preceding estimate
and Parseval's identity. The same argument is valid after
multiplication by $\psi$, proving \eqref{eq:transport-norm} and
\eqref{eq:summed-stationary-mass} for the retained arrays with the
factor $\psi^2$ on the right. The cutoffs make separation exact, not
merely asymptotic. Different old $m$ never contribute to the same
retained $j$, and at a shared $j$ the remaining supports in $u'=a$
are disjoint. The factor $\ee(ma')$ in
\eqref{eq:retained-packets} proves the asserted independence of the
new spare-coordinate modulus, also after combining packets at one
index. Linearity proves the statement about attached constants.

\emph{Fine temporal density.}
We first remove rounding before any fine-scale sampling. Smoothness and
compact support of the fixed $W$ give
\begin{equation}\label{eq:rounding-mass-removal}
 \frac1N\sum_j\int_\T
 |W(j/N-m\Delta_N,a)-W(j/N,a)|\dd a=o(1).
\end{equation}
For example, a uniform bound on $\partial_s W$ and $O(N)$ nonzero
summands give an $O(|m|/N)$ bound. The same holds with $\psi^2$ and
for any residue subset. Pushforward contracts total variation, so
\eqref{eq:summed-stationary-mass} and
\eqref{eq:rounding-mass-removal} allow us to replace the retained
squared coefficient measures by
$N^{-1}W(j/N,a)\psi(a)^2\dd a$, at total error $o(1)$, before
pushing to temporal frequency. In particular, the actual lattice
$j/N$ is never translated by $m\Delta_N$.

On the retained arc, the change of variable $y=sJ(a)$ gives
\begin{equation}\label{eq:K-jacobian}
 K(s,y)=\frac{W(s,a)\psi(a)^2}{|s|J'(a)},
 \qquad a=J|_{I'}^{-1}(y/s),
\end{equation}
with zero extension. The compact support of $\psi$ inside $I'$ makes
this extension continuous (in fact smooth away from $s=0$).
The real temporal frequency of a new angular index $j$ is
\begin{equation}\label{eq:output-frequency}
 t=j\sigma_N(a)=j/N+(j/N)J(a)/N.
\end{equation}
The density of the pushforward of
$N^{-1}W(j/N,a)\psi(a)^2\dd a$ under this map is exactly
$K(j/N,Nt-j)$. The factor $N$ in the density change of variable
cancels the coefficient mass $N^{-1}$.

Choose $Y$ so that $K(s,y)=0$ for $|y|>Y$. For fixed $t$, at most
$2Y+2$ integers contribute to either
$\sum_j K(j/N,Nt-j)$ or $\sum_jK(t,Nt-j)$; on every contributing
term $|j/N-t|\le Y/N$. Uniform continuity of $K$, and a fixed
bounded interval containing the $t$ supports for large $N$, prove
that replacing $j/N$ by $t$ costs $o(1)$ in $L^1(\R)$.
The same proof works after restriction to $j\equiv h\pmod P$.
This is \eqref{eq:fine-density-unfolded}. Folding modulo one
contracts the $L^1$ norm and finishes the proof.
\end{proof}

It is useful to emphasize the normalization in
\eqref{eq:fine-density-unfolded}. The periodization
$\sum_k K(t,y-kP)$ is a density relative to length on the circle of
length $P$, so its mean value is
\begin{equation}\label{eq:residue-mean}
 \frac1P\int_\T W(t,a)\psi(a)^2\dd a.
\end{equation}
For $P=1$, folding $t=z+l$ leaves the fast phase $Nz$ unchanged,
since $N$ is an integer. For general $P$, the folded expression has
phase $N(z+l)-h$ modulo $P$; no divisibility condition on $N$ is
imposed.

\begin{lemma}[Fast-variable averaging]\label{lem:fast-average}
Let $P>0$, and let $G(t,y)$ be continuous, compactly supported in $t$,
and $P$-periodic in $y$. Then, uniformly in real offsets $\beta_N$,
\begin{equation}\label{eq:fast-average}
 \int_\R G(t,Nt+\beta_N)\dd t
 \longrightarrow \frac1P\int_\R\int_0^P G(t,y)\dd y\dd t.
\end{equation}
In particular, suppose $w\ge0$ is continuous and compactly supported
in a compact set $K\subset\R\setminus\{0\}$. For continuous
integrable densities $\rho,\widetilde\rho$ define
\[
 \mathcal P_{s,P}\rho(y)
     =\sum_{k\in\Z}|s|^{-1}\rho((y+kP)/s).
\]
Assume these periodizations are jointly continuous for $s\in K$.
Then
\begin{align}\label{eq:law-diagonal-comparison}
 &\lim_{N\to\infty}
 \int_\R w(s)
 |\mathcal P_{s,P}\rho(Ns+\beta_N)
       -\mathcal P_{s,P}\widetilde\rho(Ns+\beta_N)|\dd s\\
 &\hspace{12mm}=
 \frac1P\int_\R w(s)\int_0^P
 |\mathcal P_{s,P}\rho(y)-\mathcal P_{s,P}\widetilde\rho(y)|
       \dd y\dd s
 \le\frac{\|w\|_1}{P}\|\rho-\widetilde\rho\|_1.\notag
\end{align}
Joint continuity is automatic, for example, if the densities are
continuous with bounds $C/(1+y^2)$, or are continuous and compactly
supported. Neither density need have mass one.
\end{lemma}

\begin{proof}
Partition the $t$ axis into intervals of length $P/N$ on each of
which $Nt+\beta_N$ traverses one period. On each interval replace
the slow argument $t$ of $G$ by an endpoint. Uniform continuity on
the fixed compact support bounds the integrated error by a constant
times its modulus of continuity at $P/N$, with an additional
$O(N^{-1})$ boundary error. The resulting averages are Riemann sums
for the right side of \eqref{eq:fast-average}. This proves the first
assertion, uniformly in the offset. Apply it to the continuous
absolute difference in \eqref{eq:law-diagonal-comparison}.
The integral over one period of that difference is bounded by
\[
 \sum_k\int_0^P |s|^{-1}
       |\rho((y+kP)/s)-\widetilde\rho((y+kP)/s)|\dd y
       =\|\rho-\widetilde\rho\|_1.
\]
Finally, on a compact set of nonzero $s$, quadratic tail bounds give
a uniform summable bound $C'/(1+k^2)$ on the periodizing summands,
for $0\le y\le P$. Uniform convergence proves joint continuity.
Compactly supported densities are handled in the same way with
only finitely many nonzero summands on that compact set.
\end{proof}

In the special case $|A(u,a)|$ is independent of $a$,
\eqref{eq:W-definition} has the form $W(s,a)=w(s)$.
If $\widetilde\rho\dd y=J_*(\psi^2\dd a)$, then
$K(s,y)=w(s)|s|^{-1}\widetilde\rho(y/s)$. Thus replacing the
realized jitter law by a regular target law $\rho$ in the
fine-density prescription costs, after taking $N$ sufficiently
large, at most
$\|w\|_1\|\rho-\widetilde\rho\|_1+o(1)$ for the full density.
One sums this estimate over the fixed packet list and folds modulo
one. This is an averaged estimate along the fast variable, not a
pointwise estimate of periodized densities from total variation.
It is applied only after \eqref{eq:rounding-mass-removal}, so the
family being averaged is fixed.

\begin{lemma}[Realization of jitter laws]\label{lem:law-approx}
For any strictly positive smooth probability density $\rho$ on
$\R$ and any $\eps>0$, there are a smooth real circle function $J$,
an open arc $I'$ with $|\T\setminus I'|<\eps$ and $J'>0$ on $I'$,
and $0\le\psi\le1$ in $C_c^\infty(I')$, such that
\[
 J_*(\psi^2\dd a)=\widetilde\rho(y)\dd y,
 \qquad \widetilde\rho\in C_c^\infty(\R),
 \qquad \|\rho-\widetilde\rho\|_1<\eps.
\]
Here the cutoff law is a subprobability measure; no renormalization
is made. Also $\int_\T(1-\psi^2)\dd a<\eps$.
\end{lemma}

\begin{proof}
Let $F(y)=\int_{-\infty}^y\rho(v)\dd v$. Strict positivity makes
$F$ a smooth diffeomorphism from $\R$ to $(0,1)$. Choose
$0<\delta<\min(1/8,\eps/4)$ and a smooth function $\xi$ compactly
supported in $(\delta/2,1-\delta/2)$ and equal to one on
$[\delta,1-\delta]$. In the circle coordinate $0<a<1$, set
$J(a)=\xi(a)F^{-1}(a)$ and extend it by zero near the identified
endpoints. It is a smooth bounded circle function, and on
$I'=(\delta,1-\delta)$ it has derivative
$1/\rho(F^{-1}(a))>0$. Choose $\psi\in C_c^\infty(I')$ with
$0\le\psi\le1$ and $\psi=1$ on $[2\delta,1-2\delta]$.
Change of variables gives
\[
 \widetilde\rho(y)=\rho(y)\psi(F(y))^2\le\rho(y),\qquad
 \|\rho-\widetilde\rho\|_1
       =1-\int_\T\psi^2\dd a\le4\delta<\eps.
\]
Its support is compact because that of $\psi$ is compactly inside
$(0,1)$. All the claims follow.
\end{proof}

\begin{lemma}[Weighted neutral refresh]\label{lem:neutral}
Given a twist, its fixed physical chart, a finite disjoint smooth
packet list with color densities $p_r$, a positive integer $P$,
and a positive tolerance, one can perform a passage with the
following simultaneous properties. Its physical map and inverse
lie in any prescribed smooth neighborhood of the old ones. Its
new good-arc complement is arbitrarily small. Every color vector
is transported within the prescribed $L^2$ tolerance, retaining
constant tags and disjointness. The new packet moduli are
independent of the new spare coordinate. The output densities
$p'_r$ and the densities $p'_{r,h}$ contributed by output indices
$j\equiv h\pmod P$ satisfy
\begin{equation}\label{eq:neutral-densities}
 \|p'_r-p_r\|_1<\eps,
 \qquad \|p'_{r,h}-p_r/P\|_1<\eps
 \quad(0\le h<P),
\end{equation}
where $\eps>0$ may be prescribed arbitrarily.
The integer $c$ may first be required to exceed any fixed bound,
and the retained supports may use any sufficiently small fixed
enlargements of \eqref{eq:stationary-image}. If colors were
disjoint in their old angular indices alone, their output angular
index sets are still disjoint. The final passage integer $N$ may
exceed any prescribed bound.
\end{lemma}

\begin{proof}
Fix the old finite data and choose $c$ as in Lemma~\ref{lem:passage},
also satisfying its prescribed lower bound. Identify all but a
small interval of the $a$ circle with an interval in $\R$.
Choose an open arc $I'$ whose complement is as small as requested,
and $0\le\psi\le1$ compactly supported there and equal to one
except on a set of sufficiently small length. The boundedness of
the finitely many old amplitudes allows us to make both their
$L^2$ cutoff errors and the removed squared masses arbitrarily
small. Fix this cutoff before proceeding.

For a positive parameter $L$, take a smooth periodic function
$J_L$ equal to $La+\beta$ on $I'$, extending it smoothly across
the complementary interval. Such an extension exists by taking
a smooth cutoff equal to one on the closure of $I'$ in a slightly
larger real coordinate interval. In particular $J_L'=L$ on $I'$.
Put $f_s(a)=W(s,a)\psi(a)^2$, extended by zero in this real
coordinate. The periodization of its pushforward by $sJ_L$ is
\begin{equation}\label{eq:neutral-riemann-sum}
 \sum_{k\in\Z}\frac1{|s|L}
 f_s\left(\frac{(y+kP)/s-\beta}{L}\right).
\end{equation}
This is $1/P$ times a Riemann sum for $\int f_s(a)\dd a$, of mesh
$P/(|s|L)$, with an arbitrary lattice offset and possibly reversed
orientation. The error is bounded by
\begin{equation}\label{eq:neutral-riemann-error}
 \frac{\|\partial_a f_s\|_{L^1(\R)}}{|s|L}.
\end{equation}
Indeed compare each lattice value times the mesh with the integral
on its cell, using the fundamental theorem of calculus and summing
the absolute derivative. Since $s$ ranges over a fixed compact set
away from zero, this bound tends to zero uniformly in $s$ and $y$.
It remains valid for arbitrary fixed dependence of the original
amplitude on $a$.

Choose $L$ large enough for these uniform errors to be as small as
needed for the finite packet list. Now $J_L$ and all its norms are
fixed. Choose the support cutoffs in Lemma~\ref{lem:passage}, and
then take $N$ sufficiently large. Its fine-density rule, together
with \eqref{eq:neutral-riemann-sum}, approximates each unfolded
residue density by $P^{-1}\int W(s,a)\psi(a)^2\dd a$ in $L^1$.
Folding modulo one gives $p_r/P$, up to the small cutoff loss;
the integer $mc$ in $mV$ has no effect on this folding. The
estimate is uniform in the periodic phase, so the folding offsets
$Nl$ cause no difficulty. Summing the residue estimates gives the
total density estimate, upon initially making each error smaller
by the fixed factor $P$ if necessary. Norm transport, exact
support properties, and smooth proximity follow from the same
passage lemma. All constraints on $N$ are lower bounds imposed
after the physical chart, the amplitudes, $c$, the cutoff, $L$,
and the support neighborhoods have been fixed.
\end{proof}

\begin{remark}\label{rem:finite-family-passages}
A common passage can handle finitely many packet lists, requiring
disjointness only within each list. Choose $c$ using the union of
their old indices, choose the neutral cutoff and slope to meet all
their finitely many mass and density bounds, and choose $N$ above
all the resulting lower bounds. Support cutoffs are chosen
separately for each list, disjoint within that list. The norm and
density conclusions then hold for every list; no orthogonality
between different lists is asserted.
\end{remark}

All these operations preserve real data. Include together the
packets $(m,A)$ and $(-m,\overline A)$, attach conjugate tags, and
choose $\chi_{-m}(s,a)=\chi_m(-s,a)$. The exact coefficients obey
$C_{-m,-j}=\overline{C_{m,j}}$, and all other cutoffs are real.
Thus conjugate packets and conjugate colors can be retained at
every passage.

\section{Small Fourier signals}
\label{sec:signals}

The purpose of this section is to modulate a packet's spectral density
according to its angular index.  The size of the modulation will depend on
a fixed band of \emph{normalized} indices, but not on the scale of those
indices.  The integer added to the old speed is allowed to depend on the
actual, finite set of indices.

Fix numbers $0<b<D_0$ and a function
\begin{equation}
 d\in C_c^\infty(\R\setminus\{0\};\C),
 \qquad d(-x)=\overline{d(x)}.
 \label{eq:signal-d}
\end{equation}
A scale is a positive real number $B$, and the normalized index of a
packet at angular index $m$ is $x=m/B$.

\begin{lemma}[Transmission of an index signal]
\label{lem:signal}
For the fixed data $b,D_0,d$ there are an integer $P\ge2$ and a number
$\gamma_0>0$ with the following property.  Any fixed
$0<\gamma\le\gamma_0$ can be used, independently of $B$.

Let $S$ be an intermediate map, and let
\[
 F_r=\sum_{m\in\mathcal M_r} A_{r,m}(u,a)\ee(m\theta)
\]
be finite smooth packet data in its good arc.  Suppose that the finite
sets $\mathcal M_r\subset\Z\setminus\{0\}$ are mutually disjoint,
$|A_{r,m}(u,a)|$ is independent of $a$, and
$b\le |m|/B\le D_0$.  Write $p_{r,m}$ for the spectral density of the
indicated packet. Given any $\eps,\delta_x,\delta_z>0$ and prescribed
positive bounds for the physical $L^2$ transport errors, a passage produces
finite output packet sums $F_{r,m}^+$, approximating the corresponding
input packets in physical $L^2$, whose spectral densities satisfy
\begin{equation}
 \sum_{r,m}\left\|p_{r,m}^+
 -p_{r,m}(z)\left[1+2\gamma\1_{\gcd(m,P)=1}
       \Rea\{d(m/B)\ee(-Nz)\}\right]\right\|_{L^1(\T)}
 <\eps.
 \label{eq:signal-conclusion}
\end{equation}
The new scale is $B^+=NcB$, and the new speed is $\sigma_N$ from
Lemma~\ref{lem:passage}. For every point $(u',a')$ in the coefficient
support of a retained output packet of index $j$ descending from
$(r,m)$, there is $(u,a)\in\supp A_{r,m}$ such that
\begin{equation}
 \left|\frac{j}{B^+}-\frac mB\right|<\delta_x,
 \qquad
 \dist_\T\bigl(j\sigma_N(u'),m\phi(u)\bigr)<\delta_z.
 \label{eq:signal-support-closeness}
\end{equation}
Its modulus is independent of the new spare coordinate. Distinct old
angular indices have disjoint output angular-index supports.
Conjugate symmetry is preserved when present.

The passage and its inverse can lie in any prescribed smooth
neighborhood of $S$ and $S^{-1}$ in physical coordinates, and the new
good arc can have arbitrarily small complement.  The choices can be made
in the order $c$, a bounded smooth jitter and cutoffs, and finally the
integer $N$.  After the preceding choices, every sufficiently large
integer $N$ works.  In particular an arbitrary lower bound on $N$ can
be imposed.
\end{lemma}

\begin{proof}
\emph{A positive law.}
We first construct a positive law without using $B$.  We use the
positive-exponent Fourier convention
\[
 \widehat q(v)=\int_\R q(R)\ee(vR)\,\dd R.
\]
Choose $0<\alpha<b/2$, put $a_0=(2\alpha)^{-1}$, and define
\begin{equation}
 s_\alpha(R)=\alpha\left(\frac{\sin(\pi\alpha R)}
                                      {\pi\alpha R}\right)^2,
 \qquad
 q_0(R)=\frac1{2a_0}\int_{-a_0}^{a_0}s_\alpha(R-t)\,\dd t,
 \label{eq:signal-baseline}
\end{equation}
with the removable value at the origin understood.  Fourier inversion
for an interval, followed by convolution, gives
\[
 \widehat{s_\alpha}(v)=(1-|v|/\alpha)_+.
\]
Consequently $q_0$ is a smooth probability density with Fourier support
in $[-\alpha,\alpha]$.  It is strictly positive: the interval of
translations in \eqref{eq:signal-baseline} cannot be contained in the
discrete zero set of $s_\alpha$.

Here a quantitative lower tail is useful.  Since the averaging interval
has length $1/\alpha$, it contains exactly one period of
$\sin^2(\pi\alpha(R-t))$ as a function of $t$.  For
$|R|\ge2a_0+1$ this gives
\[
 \frac1{2\pi^2\alpha(|R|+a_0)^2}
 \le q_0(R)
 \le\frac1{\pi^2\alpha(|R|-a_0)^2}.
\]
Positivity and continuity on the remaining compact interval therefore
give constants $c_0,C_0>0$ such that
\begin{equation}
 \frac{c_0}{1+R^2}\le q_0(R)\le\frac{C_0}{1+R^2}
 \qquad(R\in\R).
 \label{eq:signal-quadratic-tails}
\end{equation}

Choose $D>b$ such that $\supp d\subset[-D,D]$, and then an integer
$L>\max(2,D/b)$.  Let $P$ be the product of all primes at most $L$.
The strict inequality in this choice will allow a small dilation of
the Fourier arguments.  Define
\[
 g(R)=\int_\R d(v)\ee(-vR)\,\dd v,
 \qquad
 C_P(h)=\sum_{k\in(\Z/P\Z)^\times}\ee(-kh/P)
 \quad(0\le h<P).
\]
The function $g$ is Schwartz by repeated integration by parts.  The
symmetry in \eqref{eq:signal-d} makes it real, and
$\int g=d(0)=0$.  The number $C_P(h)$ is real because the set of units
is invariant under negation.  Set
\begin{equation}
 Q_h(R)=\frac1P\bigl[q_0(R)+\gamma C_P(h)g(R)\bigr].
 \label{eq:signal-joint-law}
\end{equation}
For $G=\sup_R(1+R^2)|g(R)|$ and $C_* =\max_h|C_P(h)|$, choose
$\gamma_0>0$ so small that
\[
 \gamma_0 C_*G\le c_0/2,
 \qquad 2\gamma_0\|d\|_\infty\le1/2;
\]
a zero factor imposes no restriction.  Then, for every permitted
$\gamma$,
\begin{equation}
 \frac{q_0(R)}{2P}\le Q_h(R)\le\frac{3q_0(R)}{2P},
 \qquad \int_\R Q_h(R)\,\dd R=\frac1P.
 \label{eq:signal-positive-law}
\end{equation}
Thus $Q_h(R)\,\dd R$ defines a probability law on
$\R\times\{0,\ldots,P-1\}$.  Finite Fourier inversion gives, for every
$v\in\R$ and $k\in\Z$,
\begin{equation}
 \sum_{h=0}^{P-1}\int_\R Q_h(R)\ee(vR+kh/P)\,\dd R
 =\widehat q_0(v)\1_{k\equiv0\ (P)}
       +\gamma d(v)\1_{\gcd(k,P)=1}.
 \label{eq:signal-joint-fourier}
\end{equation}
Indeed the sum over $h$ pairs $C_P(h)$ with the character at $k$,
leaving precisely the indicator that $k$ is a unit.  All choices up to
this point are independent of $B$, the map, and its finite packet list.

\emph{Moments and the finite sieve.}
For fixed $B$ and an integer $c>0$, consider the random variable
\begin{equation}
 J_*=(R/B+h/P)/c
 \label{eq:signal-ideal-jitter}
\end{equation}
under this law.  Its density is
\[
 \rho_{B,c}(y)=cB\sum_{h=0}^{P-1}Q_h\bigl(B(cy-h/P)\bigr).
\]
It is a strictly positive smooth probability density.  The bounds
\eqref{eq:signal-positive-law} imply quadratic upper and lower tails,
with constants now allowed to depend on $B,c$.

We calculate the periodized density of $sJ_*$.  Write $x=m/B$ and
$s=mct$.  For its positive integer Fourier moments, the exact formula is
\begin{equation}
 M_\ell(t)=\frac1P\sum_{h=0}^{P-1}\ee(\ell mt h/P)
 \bigl[\widehat q_0(\ell xt)
             +\gamma C_P(h)d(\ell xt)\bigr],\qquad \ell\ge1.
 \label{eq:signal-exact-moment}
\end{equation}
At $t=1$, \eqref{eq:signal-joint-fourier} applies.  Its baseline term
vanishes since $|\ell x|\ge b>\alpha$.  If $\ell\ge2$ and
$d(\ell x)\ne0$, then $\ell\le D/b<L$.  Some prime divisor of $\ell$
therefore divides $P$, so $\gcd(\ell m,P)>1$.  It follows that
\begin{equation}
 M_1(1)=\gamma d(x)\1_{\gcd(m,P)=1},
 \qquad M_\ell(1)=0\quad(\ell\ge2).
 \label{eq:signal-ideal-moments}
\end{equation}

For finite $c$ the discrete phases in
\eqref{eq:signal-exact-moment} need not be residue characters.  We
estimate their errors before using the sieve.  Choose
$\tau\in(1/2,1)$ so that $D/(\tau b)<L$, and take $c$ large enough
that
\[
 t_N(u)=1+\frac{\phi(u)+\Delta_N}{c}\ge\tau,
 \qquad |\Delta_N|\le1/N,
\]
on the old supports, uniformly for $N\ge1$.  The baseline term in
\eqref{eq:signal-exact-moment} then vanishes for \emph{every}
$\ell\ge1$, since $|\ell xt_N|\ge\tau b>\alpha$.
The perturbation also vanishes identically when
$\ell>D/(\tau b)$.  Hence nonintegral discrete phases cannot introduce
an unbounded collection of harmonics.  For the fixed remaining finite
set, put $C_1=P^{-1}\sum_h|C_P(h)|$.  The elementary bound
$|\ee(v)-\ee(w)|\le2\pi|v-w|$ gives
\begin{equation}
 |M_\ell(t)-M_\ell(1)|
 \le\gamma C_1\ell
 \bigl[D_0\|d'\|_\infty+2\pi|m|\|d\|_\infty\bigr]|t-1|.
 \label{eq:signal-finite-c}
\end{equation}
We may thus make the sum of all moment errors arbitrarily small by
choosing $c$ after the actual finite set of old indices.  In particular
$c$ is allowed to grow with $\max|m|$; $\gamma$ does not need to decrease.
The same choice can make the ranges $m(c+\phi(\T))$ for distinct old
indices disjoint and bounded away from zero.

To pass from moments to densities, let
\[
\mathcal P_s\rho(y)
   =\sum_{k\in\Z}\frac1{|s|}\rho((y+k)/s),\qquad y\in\T.
\]
Once $c$ has been chosen, all the relevant $s$ ranges are fixed compact
sets away from zero.  No bound on these sets uniform over future
passages is needed.
For $s$ in a fixed compact subset of $\R\setminus\{0\}$, the quadratic
upper tail of $\rho_{B,c}$ makes this a jointly continuous function of
$(s,y)$, with uniformly summable tails.  Its Fourier coefficients in
the positive-exponent convention are the moments just computed.
The compact bandwidth shows that its Fourier series is a finite sum
on the relevant branches.  Fourier uniqueness and
\eqref{eq:signal-ideal-moments}--\eqref{eq:signal-finite-c} therefore
give, uniformly on those branches and in $y$,
\begin{equation}
 \mathcal P_s\rho_{B,c}(y)
 =1+2\gamma\1_{\gcd(m,P)=1}\Rea\{d(m/B)\ee(-y)\}
       +o_{c\to\infty}(1).
 \label{eq:signal-periodic-density}
\end{equation}
The error here is for the fixed old list and scale $B$.
The negative exponent follows directly from our convention:
$\int\mathcal P_s\rho(y)\ee(\ell y)\,\dd y=M_\ell$.

\emph{Realization and support.}
We now realize this law by a passage.  Lemma~\ref{lem:law-approx}
provides a bounded smooth circle function $J$, strictly increasing on
a good arc, and a smooth cutoff $\psi$ compactly supported in that arc,
such that the subprobability law
\[
 \widetilde\rho=J_*(\psi(a)^2\,\dd a)
\]
is arbitrarily close to $\rho_{B,c}$ in total variation.  Its missing
mass is included in this error.  The good arc has as small a complement
as requested.  Since $|A_{r,m}|$ is independent of $a$, its stationary
weight is a function of $s$ alone:
\[
 w_{r,m}(s)=\frac{|A_{r,m}(u,a)|^2}{|m\phi'(u)|},
 \qquad m(c+\phi(u))=s,
\]
extended by zero.  For such a nonnegative compactly supported weight,
Lemma~\ref{lem:fast-average} and contraction of total variation under
pushforward and periodization give
\begin{align}
 &\lim_{N\to\infty}\int_\R w_{r,m}(s)
   |\mathcal P_s\widetilde\rho(Ns)
               -\mathcal P_s\rho_{B,c}(Ns)|\,\dd s \notag\\
 &\qquad=\int_\R w_{r,m}(s)\int_\T
   |\mathcal P_s\widetilde\rho(y)
               -\mathcal P_s\rho_{B,c}(y)|\,\dd y\,\dd s \notag\\
 &\qquad\le\|w_{r,m}\|_1\|\widetilde\rho-\rho_{B,c}\|_{L^1(\R)}.
 \label{eq:signal-law-comparison}
\end{align}
Here densities of subprobabilities are used in the $L^1$ norm.
The realized law has a continuous compactly supported density; the
ideal law has the uniform tail bound established above.  These facts
supply the continuity and domination required for fast averaging.
The independence of the weight from $a$ is essential in this step.

Lemma~\ref{lem:passage} identifies the output density, up to an
arbitrarily small $L^1$ error, as the folding to $\T$ of
$w_{r,m}(s)\mathcal P_s\widetilde\rho(Ns)$.
In using that lemma the rounding correction is removed at the
coefficient-mass level: its replacement of $w_{r,m}(s)$ by
$w_{r,m}(s-m\Delta_N)$ costs $o(1)$ in the summed masses.  It does
not translate the fast argument $Ns$.  Alternatively,
\eqref{eq:signal-finite-c} also holds uniformly with $t=t_N$.
Since $cm$ and $N$ are integers,
\[
 p_{r,m}(z)=\sum_{k\in\Z}w_{r,m}(z+k),
 \qquad \ee(-N(z+k))=\ee(-Nz).
\]
Now \eqref{eq:signal-periodic-density},
\eqref{eq:signal-law-comparison}, and the passage density formula
prove \eqref{eq:signal-conclusion}.  Choose $c$ first to make its
error sufficiently small, then the law approximation, and then $N$.
The finite packet list allows all errors to be summed.

For completeness we give the simultaneous support estimates, which
will also apply to neutral passages.  Put
$M=\max_{r,m}|m|$, with $M=0$ for an empty packet list.
By the support cutoff in
Lemma~\ref{lem:passage}, at each retained point of an output packet at
index $j$ there is a point $u$ in its old support for which
\begin{equation}
 j/N=m(c+\phi(u)+\Delta_N)+\eta,
 \qquad |\eta|\le\eta_0,
 \label{eq:signal-retained-support}
\end{equation}
where the image enlargement $\eta_0$ can be as small as desired.
If $x'=j/(NcB)$ and $z'=j/N+(j/N)J(a)/N$ is the new temporal
frequency, then
\begin{align}
 |x'-m/B|
 &\le\frac{D_0(\|\phi\|_\infty+1)}c+\frac{\eta_0}{cB},
 \label{eq:signal-index-drift}\\
 \dist_\T(z',m\phi(u))
 &\le\frac MN+\eta_0
  +\frac{[M(c+\|\phi\|_\infty+1)+\eta_0]\|J\|_\infty}{N}.
 \label{eq:signal-temporal-drift}
\end{align}
Choose $c$ for the first bound and index separation, then $\eta_0$
and the bounded $J$, and finally $N$ for the second bound.  A large
bound on $J$ only increases the eventual requirement on $N$.
These choices also give physical smooth closeness of the passage and
its inverse and the required $L^2$ approximations, by
Lemma~\ref{lem:passage}.

Every new coefficient descending from old index $m$ has the form
$\ee(ma_{\rm new})$ times a scalar function of the new speed
coordinate.  Because different old indices have disjoint retained
new index ranges, their spare-coordinate phases are never added at
one new index.  Thus the resulting packet moduli are independent of
the new spare coordinate.  All cutoffs can be chosen in conjugate
pairs, so conjugation is retained as well.
\end{proof}

We next remove the coprimality restriction by placing a neutral passage
before transmission.  The output residue classes of the neutral
passage, rather than the input residue classes, are the ones tested in
the transmission.

\begin{lemma}[Refresh and transmission]
\label{lem:refresh-transmit}
Fix $b,D_0,d,P,\gamma$ as in Lemma~\ref{lem:signal}, and put
\begin{equation}
 \lambda=\gamma\frac{|(\Z/P\Z)^\times|}{P}>0.
 \label{eq:signal-lambda}
\end{equation}
Suppose finite colored packet data have mutually disjoint angular-index
sets and moduli independent of the spare coordinate.  Suppose their
normalized indices lie in a fixed compact subset of
$\{x:b<|x|<D_0\}$.  Let $p_r$ be the color densities, and let
$d_r:\T\to\C$ be continuous functions such that on the support of
every input packet of color $r$,
\begin{equation}
 |d(m/B)-d_r(m\phi(u)\bmod1)|\le\omega.
 \label{eq:signal-color-match}
\end{equation}
Then a neutral passage followed by a transmission gives output colors
with
\begin{equation}
 \sum_r\left\|p_r^+
  -p_r(z)\bigl[1+2\lambda\Rea\{d_r(z)\ee(-Nz)\}\bigr]
                 \right\|_1
 \le 2\gamma\omega\sum_r\|p_r\|_1+\eps,
 \label{eq:signal-refreshed}
\end{equation}
where $\eps>0$ is arbitrary and $N$ is the second passage integer.
The physical map and inverse, the physical packet vectors, and the
normalized-index and temporal-frequency displacements can all be
controlled with arbitrary prescribed positive error bounds.  Output
packet moduli remain independent of the spare coordinate, and colors
remain disjoint in angular indices.  The final $N$ can exceed any
prescribed bound.  Conjugate-symmetric data can be treated
conjugate-symmetrically.

If different colors start in separated compact normalized-index
sets, any fixed sufficiently small enlargements of those sets can
contain their outputs.  The same conclusion holds through any fixed
finite number of such pairs, with positive margins around the band
$b<|x|<D_0$.  If the input supports satisfy
$\dist_\T(m/B,m\phi(u))\le\kappa$, the output supports satisfy the
same inequality with $\kappa$ increased by an arbitrarily small amount.
\end{lemma}

\begin{proof}
Use Lemma~\ref{lem:neutral} with residue modulus $P$.  Denote the
refreshed scale by $\widetilde B$, its packet densities by
$\widetilde p_{r,m}$, and its color-residue densities by
\[
 \widetilde p_{r,h}=\sum_{m\equiv h\ (P)}\widetilde p_{r,m}.
\]
That lemma makes
\begin{equation}
 E_r=\sum_{h=0}^{P-1}
       \|\widetilde p_{r,h}-p_r/P\|_1
 \label{eq:signal-residue-error}
\end{equation}
arbitrarily small.  It also retains disjoint output angular-index
sets and gives moduli independent of the new spare coordinate.
To control its support, first choose its $c$ as large as necessary
for \eqref{eq:signal-index-drift}; next choose its cutoff and broad
linear jitter for the residue equalization; finally choose its passage
integer for \eqref{eq:signal-temporal-drift}.  Thus equalizing residues
does not prevent either displacement from being arbitrarily small.

For a continuous function $f$ on its relevant compact domain, denote
its modulus of continuity by
$\nu_f(t)=\sup_{\dist(v,w)\le t}|f(v)-f(w)|$.
If the refresh changes normalized indices by at most $\delta_x$ and
temporal frequencies by at most $\delta_z$, its retained color
supports satisfy
\begin{equation}
 |d(m/\widetilde B)-d_r(z)|
 \le\rho_r:=\omega+\nu_d(\delta_x)+\nu_{d_r}(\delta_z).
 \label{eq:signal-refreshed-match}
\end{equation}
This follows pointwise by linking an output support point to an old
one and using \eqref{eq:signal-color-match}.  In particular the extra
terms tend to zero.  The compact margin in the normalized band lets
us apply Lemma~\ref{lem:signal} with the same $P,\gamma$ after this
refresh, although its actual indices and scale may be very large.

By that lemma the transmission density, with an arbitrarily small
summed error $T_r$ for color $r$, is
\[
 \widetilde p_r(z)+2\gamma\Rea\left\{\ee(-Nz)
    \sum_{\gcd(m,P)=1}d(m/\widetilde B)
                                   \widetilde p_{r,m}(z)\right\}.
\]
Replacing the coefficient in this sum by $d_r(z)$ costs at most
$2\gamma\rho_r\|\widetilde p_r\|_1$ in $L^1$.  Also,
\[
 \|\widetilde p_r-p_r\|_1\le E_r,
 \qquad
 \left\|\sum_{h\in(\Z/P\Z)^\times}\widetilde p_{r,h}
          -\frac{|(\Z/P\Z)^\times|}{P}p_r\right\|_1\le E_r.
\]
Consequently the color's total error against the target in
\eqref{eq:signal-refreshed} is bounded by
\begin{equation}
 T_r+(1+2\gamma\|d_r\|_\infty)E_r
       +2\gamma\rho_r(\|p_r\|_1+E_r).
 \label{eq:signal-refresh-error}
\end{equation}
Choose the refresh errors, displacements, and transmission errors
small enough and sum this bound over the finite list.  This proves
\eqref{eq:signal-refreshed}, including the exact value of $\lambda$.
Only residue balancing after summing over a color was needed, because
\eqref{eq:signal-refreshed-match} holds pointwise before the sum.

The remaining assertions follow from the two passage lemmas and
\eqref{eq:signal-index-drift}--\eqref{eq:signal-temporal-drift}.
Specifically, the triangle inequality gives
\[
 \dist_\T(x',z')\le\dist_\T(x,z)+|x'-x|+\dist_\T(z',z).
\]
Positive separation distances and the fixed finite number of passages
allow displacement budgets smaller than all band margins and color
gaps.  Constants such as $\lambda$ are fixed before these error budgets
are imposed.  The required $c$ and passage integers may then be chosen
afresh after each actual finite packet list is known.
\end{proof}

\section{Finite blocks and prediction}
\label{sec:blocks}

The signal from Lemma~\ref{lem:refresh-transmit} can be arbitrarily
weak. We use many such signals, holding one centering function fixed
throughout each block. For a block with $n\lambda^2$ comparable to
$\delta$, the square-root integral loses $O(\delta^2)$, while the
estimator has mean-square error $O(1/\delta)$. Averaging the estimators
from $K$ blocks gives mean-square error $O((K\delta)^{-1})$ with total
loss $O(K\delta^2)$. Thus $K\delta$ can grow while $K\delta^2$ tends
to zero.

\begin{proposition}[Prediction of a tagged packet sum]
\label{prop:prediction}
Let $S$ be an intermediate map, expressed in a fixed physical
volume-preserving coordinate chart.  Let
\[
 F_r=\sum_{m\in\mathcal M_r} A_{r,m}(u,a)\ee(m\theta),
 \qquad F=\sum_r F_r,
\]
be a finite list of colors of smooth packets supported compactly in the
good arc.  Suppose that the finite sets
$\mathcal M_r\subset\Z\setminus\{0\}$ are pairwise disjoint, that
$|A_{r,m}(u,a)|$ is independent of $a$, and that the list is invariant
under complex conjugation: $F_{r^*}=\overline{F_r}$.  In particular, $F$
is real.  Let $p_r$ be the corresponding spectral densities, put
$p=\sum_r p_r$, and suppose that
\[
 \int_\T p(z)\,\dd z\le \frac32.
\]
Assign constants $b_r\in\C$ satisfying
$b_{r^*}=\overline{b_r}$.

Given $\eps,\tau,\kappa>0$ and any neighborhood of the physical map
and its inverse in the smooth topology, there are a finite sequence of
passages, a final intermediate map $S'$, final packets $F'_r$, and a
smooth function $\Psi:\T\to\C$, such that
\begin{align}
 \sum_r\|F'_r-F_r\|_2&<\eps,
 \label{eq:prediction-transport}\\
 \int_\T\sqrt{p'(z)}\,\dd z
 &\ge \int_\T\sqrt{p(z)}\,\dd z-\tau,
 \qquad p'=\sum_r p'_r,
 \label{eq:prediction-hellinger}\\
 \left\|\sum_r b_rF'_r-
       \Psi(U_{S'})\sum_r F'_r\right\|_2&<\eps.
 \label{eq:prediction-error}
\end{align}
Here the differences in \eqref{eq:prediction-transport} are measured
in the original physical $L^2$ space.  All the physical maps and their
inverses can be kept in the prescribed neighborhood.  The final good
arc $I'$ satisfies $|\T\setminus I'|<\kappa$.
The output colors again have disjoint nonzero angular-index sets,
moduli independent of the spare coordinate, and conjugation symmetry.
The multiplier can be chosen with
$\Psi(-z)=\overline{\Psi(z)}$.
\end{proposition}

We first record the precise averaging assertion needed in the proof.
In the following lemma, a Laurent polynomial in $w_1,\ldots,w_n$ means
a finite linear combination of the monomials
$w^\nu=\prod_i w_i^{\nu_i}$, with $\nu\in\Z^n$.

\begin{lemma}[Finite phase tests]
\label{lem:phase-tests}
Fix finitely many expressions
\[
 P_j(z,w)=\sum_{\nu\in\mathcal E_j}c_{j,\nu}(z)w^\nu,
 \qquad c_{j,\nu}\in L^1(\T),
\]
where each $\mathcal E_j$ is finite.  For every $\eta>0$, successively
large positive integers $L_1,\ldots,L_n$ can be chosen so that,
simultaneously for all $j$,
\begin{equation}
 \left|\int_\T P_j\bigl(z,\ee(L_1z),\ldots,\ee(L_nz)\bigr)\,\dd z
 -\int_\T\int_{\T^n}P_j\bigl(z,\ee(y_1),\ldots,\ee(y_n)\bigr)
             \,\dd y\,\dd z\right|<\eta.
 \label{eq:phase-tests}
\end{equation}
At each selection, $L_i$ may additionally be required to exceed any
finite lower bound determined by the preceding choices and the
then-current construction data.
\end{lemma}

\begin{proof}
Let $D$ bound the absolute values of all the exponents in all the
$\mathcal E_j$.  The integral of a nonconstant monomial on the product
torus is zero.  On the one-dimensional diagonal its contribution is a
Fourier coefficient of $c_{j,\nu}$ at
$-\sum_i\nu_iL_i$.
The Fourier coefficients of an $L^1$ function tend to zero: this
follows by approximation in $L^1$ by trigonometric polynomials, whose
sufficiently high Fourier coefficients vanish.  Since the coefficient
list is finite, choose $R$ so that the sum of the bounds for all these
nonconstant contributions is less than $\eta$ whenever their
frequencies have absolute value greater than $R$.
Choose successively
\[
 L_j>R+D\sum_{i<j}L_i,
\]
as well as the additional finite lower bounds.  If $j$ is the last
nonzero position of $\nu$, then
\[
 \left|\sum_i\nu_iL_i\right|
 \ge L_j-D\sum_{i<j}L_i>R.
\]
This proves the assertion.  No independence assertion about the
distribution of the entire diagonal is made or needed.
\end{proof}

\begin{proof}[Proof of Proposition~\ref{prop:prediction}]
Zero colors may be omitted.  If all colors vanish, a neutral passage
with zero output packets and $\Psi=0$ proves the assertion, including
the requested new good arc.  Hence assume that a nonzero color is
present.
Dividing all tags by $\max(1,\max_r|b_r|)$ and reducing the requested
prediction tolerance reduces the proof to $|b_r|\le1$.  The transport
and square-root tolerances remain independently adjustable.  We work
under this normalization until the end.

\paragraph{The data fixed at the beginning of a block.}
For the moment suppose that the current packet mass is at most $2$.
Write its actual densities as $p_r$, put $p=\sum_rp_r$, and on
$\{p>0\}$ set
\[
 t_r=\frac{p_r}{p},\qquad
 \beta=\sum_r t_rb_r.
\]
Set $\beta=0$ on $\{p=0\}$.  Then $|\beta|\le1$ and
$\beta(-z)=\overline{\beta(z)}$.  Fix $0<\delta\le1/100$.
There is a smooth circle function $q$ such that
\begin{equation}
 |q|\le2,\qquad q(-z)=\overline{q(z)},\qquad
 \int_\T\sqrt p\,|\beta-q|^2\,\dd z<\delta.
 \label{eq:block-centering}
\end{equation}
Indeed, convolve $\beta$ with an even nonnegative smooth approximate
identity.  The convolutions have absolute value at most $1$, have
the required symmetry, and converge almost everywhere to $\beta$.
Since $\sqrt p\in L^1$, dominated convergence gives
\eqref{eq:block-centering}.  In particular this choice does not require
any uniform bound on derivatives of $q$.

Let $A$ be any fixed bounded function on the temporal circle.  In the
iteration, $A$ will be the sum of the estimators from $k$ already
completed blocks.  Both $q$ and $A$ are fixed before any frequencies
for the present block are chosen.  The densities $p_r$ in every model
below are these actual \emph{pre-block} densities.

\paragraph{Anchoring, with the number of passages fixed first.}
Choose an integer $c_0$ large enough that the full compact ranges
$m(c_0+\phi(\T))$ for distinct old indices are separated and avoid
zero, as required by Lemma~\ref{lem:passage}.  In particular this
separates the compact packet images
\[
 m(c_0+\phi(u)),\qquad
 u\in\operatorname{proj}_u\supp A_{r,m},
\]
for all distinct old indices $m$.  Let $K_r$ be the union of the
images belonging to color $r$.
The $K_r$ are disjoint compact sets and
$K_{r^*}=-K_r$.  Choose slightly larger compact neighborhoods of them,
still disjoint and off zero, and then disjoint open neighborhoods with
compact closures.  We will retain enough room between these two
families of neighborhoods for all the support movements in the block.
There is a smooth compactly supported function on $\R\setminus\{0\}$
such that
\begin{equation}
 d(-x)=\overline{d(x)},\qquad
 d(x)=b_r-q(x\bmod1)
 \quad\hbox{on a neighborhood of the enlarged }K_r.
 \label{eq:block-band-function}
\end{equation}
For example, use disjoint smooth cutoffs equal to one near these sets,
chosen compatibly with reflection, and multiply the $r$th cutoff by
$b_r-q(x\bmod1)$.  Fix $0<b_0<D_0$ such that all the retained
neighborhoods lie strictly inside $\{b_0<|x|<D_0\}$.

Apply Lemma~\ref{lem:signal} to this fixed $d$ and these band bounds.
It gives a modulus $P$ and an arbitrarily small positive strength
$\gamma$, independently of the future angular scale.  Put
\[
 \lambda=\gamma\frac{|(\Z/P\Z)^\times|}{P}.
\]
Take $\gamma$ so small that $6\lambda<1$ and $\lambda^2\le\delta/2$,
and choose an integer $n\ge1$ with
\begin{equation}
 \frac\delta2\le n\lambda^2\le\delta.
 \label{eq:block-budget}
\end{equation}
One may take $n=\lfloor\delta/\lambda^2\rfloor$.
At this point $q,c_0,K_r,d,P,\gamma,\lambda,n$ are all fixed.
In particular they have been chosen \emph{before} the integer of the
anchor passage.  The error targets used below depend only on these
data, the already fixed $A,k$, and the requested tolerance $\nu$.
They can therefore be fixed now, before executing the anchor.

Now make a neutral passage with the already fixed $c_0$, choosing its
jitter and its integer $N_0$ for the required accuracy.  Use
$B=N_0$ as the output angular scale.  Lemmas~\ref{lem:passage}
and~\ref{lem:neutral} give color densities as close to $p_r$ in
$L^1$ as desired.  On their retained supports the normalized indices
$x=j/N_0$ lie in arbitrarily small enlargements of $K_r$.  If $J_0$
is the bounded jitter chosen for this passage, their actual temporal
frequencies satisfy
\[
 z=x+xJ_0(a)/N_0\pmod1.
\]
Thus $\dist_{\T}(x\bmod1,z)$ can be made arbitrarily small after
$J_0$ is fixed, by increasing $N_0$.

We follow this anchor by $n$ refresh--transmit pairs from
Lemma~\ref{lem:refresh-transmit}. At every passage the scale changes by
$B'=NcB$. The support estimates
\eqref{eq:signal-index-drift}--\eqref{eq:signal-temporal-drift}, which
apply to neutral as well as transmitting passages, control normalized-index
and temporal-frequency movement separately. Choose $c$ and sufficiently
small support enlargements to meet the movement budgets; in particular,
the enlargement is smaller than the temporal tolerance because its
contribution there is not divided by $N$. Fix the bounded jitter $J$
next, and choose $N$ last for the remaining terms. Thus an arbitrarily
broad neutral jitter remains compatible with temporal tracking.

There are only $2n+1$ passages in this block.  Apportion small enough
movement budgets among them so that all normalized indices remain
in the neighborhoods on which \eqref{eq:block-band-function} holds,
and so that $\dist_{\T}(x\bmod1,z)$ remains below any prescribed
positive tolerance.  Uniform continuity of the fixed $q$ then makes
\begin{equation}
 |d(x)-(b_r-q(z))|
 \label{eq:block-substitution}
\end{equation}
uniformly as small as needed on every retained support of color $r$.
The positive gaps preserve disjoint angular-index supports.  Each
passage preserves the independent-of-spare-coordinate modulus
property and conjugation symmetry.

\paragraph{The product model and its approximation.}
Put $a_r(z)=b_r-q(z)$, so $|a_r|\le3$.  Denote the integer of the
$i$th transmission by $L_i$.  The model after $i$ transmissions is
\begin{equation}
 h_r^{(i)}(z)=p_r(z)\prod_{j=1}^i f_{r,j}(z),\qquad
 f_{r,j}(z)=1+2\lambda\Rea\{a_r(z)\ee(-L_jz)\}.
 \label{eq:block-model}
\end{equation}
All factors are positive, and $f_{r,j}\le M=1+6\lambda<2$.
Let $g_r^{(i)}$ be the actual densities after these pairs, with
$g_r^{(0)}$ the anchor densities.  The refresh--transmit lemma,
including the substitution \eqref{eq:block-substitution}, allows
\[
 \sum_r\|g_r^{(i)}-g_r^{(i-1)}f_{r,i}\|_1<\eta_i,
 \qquad
 \sum_r\|g_r^{(0)}-p_r\|_1<\eta_0.
\]
Indeed the substitution costs at most
$2\gamma\omega\sum_r\int g_r^{(i-1)}$ if its uniform error is
$\omega$; the remaining error is freely adjustable.  We may first
choose the movement tolerance small enough for all the finitely many
$\eta_i$, and then implement each pair on its actual finite input.
Consequently
\begin{equation}
 E_i:=\sum_r\|g_r^{(i)}-h_r^{(i)}\|_1
 \le ME_{i-1}+\eta_i,
 \qquad
 E_n\le M^n\eta_0+\sum_{i=1}^nM^{n-i}\eta_i.
 \label{eq:block-error-propagation}
\end{equation}
For any prescribed $e_0>0$, it suffices to choose every
$\eta_i<e_0/((n+1)M^n)$ to obtain $E_n<e_0$.
All quantities in this accuracy requirement were fixed before the
anchor.  Even if $\lambda$ is extremely small and $n$ and $M^n$ are
extremely large, they are finite known numbers.

The frequencies $L_i$ remain free to exceed arbitrary successive
lower bounds.  In particular, Lemma~\ref{lem:phase-tests} can be
applied to any finite list of Laurent polynomials whose coefficient
functions depend only on the pre-block data
\begin{equation}
 p_r,\ t_r,\ q,\ A,\ b_r,\ n,\ \lambda.
 \label{eq:block-fixed-data}
\end{equation}
The actual packets depend on previous choices, but appear only in
\eqref{eq:block-error-propagation}, not in this fixed coefficient
list.  At the $i$th pair, first choose its refresh and all its
transmission data preceding the passage integer; then choose $L_i$
large enough for both the physical construction and the phase-test
lower bound.  Thus the averaging and the construction impose
compatible lower bounds on the same integer.

\paragraph{The square-root integral.}
Introduce independent Haar phases $Z_1,\ldots,Z_n$ solely for
computing the product-torus averages in
Lemma~\ref{lem:phase-tests}.  At each fixed $z$ put
\[
 H_r=\prod_{i=1}^n
       [1+\lambda(a_r\overline{Z_i}+\overline{a_r}Z_i)],
 \qquad R=\sum_r t_rH_r.
\]
The following pointwise calculations are on $\{p>0\}$; all their
integrands are multiplied by $p_r$ or $\sqrt p$ and are extended by
zero on its complement.  We have
\begin{equation}
 \mathbb E H_r=1,
 \qquad
 \mathbb E(H_rH_s)
       =(1+2\lambda^2\Rea(a_r\overline{a_s}))^n,
 \qquad \mathbb E R=1.
 \label{eq:block-second-moment-identity}
\end{equation}
For real $y$ and integer $n\ge1$, the binomial expansion gives
\[
 |(1+y)^n-1-ny|
 \le\sum_{j=2}^n\frac{(n|y|)^j}{j!}
 \le\frac{(n|y|)^2}{2}\exp(n|y|).
\]
Apply this with $y=2\lambda^2\Rea(a_r\overline{a_s})$;
then $n|y|\le18\delta$.  Since $\sum_rt_r=1$, the sum of the
linear terms in \eqref{eq:block-second-moment-identity} is
$2n\lambda^2|\sum_rt_ra_r|^2$.  Therefore
\begin{equation}
 \mathbb E(R-1)^2
 \le2\delta|\beta-q|^2+162e^{18\delta}\delta^2.
 \label{eq:block-second-moment-bound}
\end{equation}
This bound is uniform in the number of colors and in all the later
geometric parameters.

For every $x\ge0$,
\begin{equation}
 \sqrt x\ge1+\frac{x-1}{2}-(x-1)^2.
 \label{eq:block-square-root-bound}
\end{equation}
In fact, on setting $t=\sqrt x$, the difference between the two sides
equals $(t-1)^2(t^2+2t+1/2)$.  Multiplying
\eqref{eq:block-square-root-bound} by $\sqrt p$, taking the
independent-phase expectation, and using
\eqref{eq:block-centering} and \eqref{eq:block-second-moment-bound}
shows that the averaged polynomial lower bound is at least
\[
 \int_\T\sqrt p\,\dd z
 -2\delta^2-162e^{18\delta}\delta^2\int_\T\sqrt p\,\dd z
 \ge \int_\T\sqrt p\,\dd z-300\delta^2.
\]
Here $\int\sqrt p\le\sqrt2$ and $\delta\le1/100$ were used.

The expression
\[
 \sqrt p\,[1+(R-1)/2-(R-1)^2]
\]
is a Laurent polynomial in the phases with fixed integrable
coefficients.  Lemma~\ref{lem:phase-tests} therefore transfers this
lower bound, up to any prescribed error, to the diagonal phases
$Z_i=\ee(L_i z)$.  On that diagonal
$pR=\sum_r h_r^{(n)}$.
There is no need to apply the lemma to $\sqrt R$ itself.
Finally, for nonnegative densities $g,h$ on the unit circle,
\begin{equation}
 \left|\int_\T\sqrt g-\int_\T\sqrt h\right|
 \le\left(\int_\T|g-h|\right)^{1/2},
 \label{eq:block-density-hellinger-error}
\end{equation}
by $(\sqrt g-\sqrt h)^2\le|g-h|$ and Cauchy--Schwarz.
Thus the actual output density $p^{\rm out}=\sum_rg_r^{(n)}$
satisfies, for any chosen $\nu>0$ and sufficiently accurate
construction,
\begin{equation}
 \int_\T\sqrt{p^{\rm out}}\,\dd z
 \ge\int_\T\sqrt p\,\dd z-300\delta^2-\nu.
 \label{eq:block-hellinger-step}
\end{equation}

\paragraph{Prediction and a previous estimator.}
With the same formal phases, put
\[
 D=q+\frac1{n\lambda}\sum_{i=1}^n Z_i.
\]
For a fixed color $r$ and fixed $z$, $H_r$ is a positive probability
density on the product phase torus, and under this density the phases
remain independent.  Direct integration of one factor gives
\[
 \mathbb E_{H_r}Z_i=\lambda a_r,
 \qquad
 \mathbb E_{H_r}|Z_i-\lambda a_r|^2
       =1-\lambda^2|a_r|^2.
\]
Consequently
\begin{equation}
 \mathbb E_{H_r}D=b_r,
 \qquad
 \mathbb E_{H_r}|D-b_r|^2
 =\frac{1-\lambda^2|a_r|^2}{n\lambda^2}
 \le\frac2\delta.
 \label{eq:block-estimator-moments}
\end{equation}
For the arbitrary fixed bounded previous estimator $A$, this means
that the cross term against $A-kb_r$ vanishes \emph{at each $z$}:
\[
 \mathbb E_{H_r}|A+D-(k+1)b_r|^2
 =|A-kb_r|^2+\mathbb E_{H_r}|D-b_r|^2.
\]
After multiplying by $p_r$, summing, and integrating, the increase
in this mean square is at most $4/\delta$.

The expression
\[
 \sum_r p_rH_r|A+q+(n\lambda)^{-1}\textstyle\sum_iZ_i-(k+1)b_r|^2
\]
is another finite Laurent polynomial with the fixed integrable
coefficients in \eqref{eq:block-fixed-data}.  Apply
Lemma~\ref{lem:phase-tests} to it simultaneously with the preceding
square-root lower-bound test.  Write
\begin{equation}
 D(z)=q(z)+\frac1{n\lambda}\sum_{i=1}^n\ee(L_i z)
 \label{eq:block-diagonal-estimator}
\end{equation}
for the resulting actual multiplier.  Its norm is bounded by
$2+1/\lambda$, independently of the choices of $L_i$.  In passing
from model to actual color densities, the error in the prediction
integral is at most $W E_n$, where
\begin{equation}
 W=(\|A\|_\infty+2+1/\lambda+k+1)^2.
 \label{eq:block-estimator-error-weight}
\end{equation}
This finite number is known before choosing the anchor accuracy or
any transmission frequencies.  For example, choose
\[
 e_0<\min\{(\nu/4)^2,\nu/(4W)\},
\]
and make the two phase-test errors less than $\nu/2$.
Together with \eqref{eq:block-error-propagation} this proves
\begin{equation}
 \sum_r\int_\T p_r^{\rm out}|A+D-(k+1)b_r|^2\,\dd z
 \le\sum_r\int_\T p_r|A-kb_r|^2\,\dd z
            +\frac4\delta+\nu,
 \label{eq:block-prediction-step}
\end{equation}
simultaneously with \eqref{eq:block-hellinger-step}.
The same $e_0$ and support accuracies can be imposed before the
anchor.  Thus a small strength entails a more demanding finite
accuracy requirement, without a circular parameter choice.

\paragraph{Iteration of blocks.}
All color vectors can be transported with an arbitrarily small sum
of physical norm errors at every passage.  Reserve a total such
budget smaller than
\[
 \min\{\eps,\sqrt2-\sqrt{3/2}\}.
\]
Since the sum of the color errors bounds the error of their sum,
all intermediate packet sums then have squared norm less than $2$.
This supplies the mass bound used above at every step, with room to
spare.  Smooth map and inverse errors can likewise be apportioned
among the finitely many blocks and, after $n$ is chosen in a block,
among its $2n+1$ passages.  Choosing a basic neighborhood specified
by finitely many smooth norms inside the requested neighborhood
makes this allocation explicit.  Every passage may also make the
new excluded arc arbitrarily small; apply this in particular at the
last passage.

Set $\nu\le\delta^2$ in each block.  Start with $A_0=0$, and let
$A_{k+1}=A_k+D_{k+1}$ after block $k+1$.  At the beginning of each
new block choose a \emph{new} $q$ from that block's actual input
densities.  That $q$ is then held fixed throughout its block.
The band function, sieve modulus, strength, and length may all
change from block to block.  The estimates depend only on
$|b_r|\le1$, $|q|\le2$, the mass bound $2$, and
\eqref{eq:block-budget}; they require no uniform lower bound for
the strengths and no uniform derivative bound for the successive
centering functions.

After $K$ blocks, telescoping
\eqref{eq:block-hellinger-step} and
\eqref{eq:block-prediction-step} gives
\begin{align}
 \int_\T\sqrt{p^{(K)}}\,\dd z
 &\ge\int_\T\sqrt{p^{(0)}}\,\dd z-400K\delta^2,
 \label{eq:block-total-hellinger}\\
 \sum_r\int_\T p_r^{(K)}
       |A_K/K-b_r|^2\,\dd z
 &\le\frac5{K\delta}.
 \label{eq:block-total-prediction}
\end{align}
For the second estimate the initial error is zero, and
$4/\delta+\delta^2\le5/\delta$.
Given the desired prediction tolerance $\eps$, first choose
$\delta>0$ sufficiently small and then put
\[
 K=\left\lceil\frac{10}{\eps^2\delta}\right\rceil.
\]
We can require simultaneously $\delta\le1/100$ and
\[
 400\left(\frac{10\delta}{\eps^2}+\delta^2\right)<\tau.
\]
Then $K\delta$ is large enough for
\eqref{eq:block-total-prediction}, and $K\delta^2$ is small enough
for \eqref{eq:block-total-hellinger}.  These choices precede all
blocks and do not depend on their future strengths or lengths.

Set $\Psi=A_K/K$.  Disjoint angular-index supports make all
different color cyclic subspaces orthogonal, so the spectral
density formula gives exactly
\[
 \left\|\sum_r b_rF_r^{(K)}-
          \Psi(U_{S^{(K)}})\sum_rF_r^{(K)}\right\|_2^2
 =\sum_r\int_\T p_r^{(K)}|b_r-\Psi|^2\,\dd z.
\]
This proves the required norm estimate.  Every $D$ in
\eqref{eq:block-diagonal-estimator} is smooth and has conjugation
symmetry, hence so does $\Psi$.  Restoring the original scale of
the tags proves the proposition in its stated form.
\end{proof}

\section{Preparation, whitening, and the smooth limit}
\label{sec:completion}

The finite prediction result encodes a target observable in a packet
sum while retaining nearly flat total spectral density. We prepare
arbitrary observables for that result, whiten the resulting vector to
exactly flat density, and diagonalize the stages to obtain a smooth map
whose one orbit is complete.

All vector comparisons in this section take place in the fixed physical
space $L^2(\T^3)$, even when intermediate coordinates change.  A neighborhood
of a diffeomorphism will mean a neighborhood of the pair consisting of the
map and its inverse in the $C^\infty$ topology.  For a unitary $U$, write
\[
 \calC(U,f)=\overline{\operatorname{span}\{U^nf:n\in\Z\}}.
\]
We use the packet spectral calculus from Section~\ref{sec:passages}:
for a measurable function $\alpha$ of the additive spectral variable $z$, the
coefficient of $\alpha(U_S)F$ at angular index $m$ is
$\alpha(m\phi(u)\bmod1)A_m(u,a)$.  The notation is used only when the resulting
vector has finite $L^2$ norm.

\begin{lemma}[A flat vector]\label{lem:flat-vector}
For every intermediate twist and every good arc $I$, there is a real
centered vector $v$ of spectral density $1$, using two nonzero angular
indices and coefficients supported in a compact subinterval of $I$.
\end{lemma}
\begin{proof}
Choose a positive integer $m$ and $u_0<u_1$ in $I$ such that
$m(\phi(u_1)-\phi(u_0))=1$.  This is possible by first choosing a compact
interval on which $\phi$ increases and then making $m$ large.  Put
\[
 A(u)=\1_{[u_0,u_1]}(u)
       \sqrt{\frac{|m\phi'(u)|}{2}},
 \qquad
 v(\theta,u,a)=A(u)\ee(m\theta)+A(u)\ee(-m\theta).
\]
The pushforward of $|A(u)|^2\dd u\,\dd a$ by each of $\pm m\phi(u)$
modulo one is $\tfrac12\dd z$.  The two angular modes are orthogonal,
so the spectral density of $v$ is $1$.  The vector is real and centered.
Its endpoint discontinuities are harmless: no smoothness is required of
the vector.
\end{proof}

We call a pair $(S,f)$ admissible if $S$ is an intermediate twist in
smooth volume-preserving coordinates, with good arc $I$, and $f$ is real,
has spectral density $1$, and has finitely many nonzero angular modes
with $L^2$ coefficients supported in a compact subset of $I\times\T$.
In particular, $f$ is centered and has norm one.

\begin{proposition}[One stage]\label{prop:stage}
Let $(S,f)$ be admissible, let $h$ be a real smooth function in physical
coordinates with $\int h=0$ and $\|h\|_\infty\le1$, and let
$\eta,\rho,\tau>0$.  In every prescribed neighborhood of $S$ there is an
admissible pair $(S',f')$, with good arc $I'$, and a finite Laurent
polynomial $P$ such that
\begin{equation}\label{eq:stage-conclusion}
 \|f'-f\|_2<\tau,\qquad
 |\T\setminus I'|<\rho,\qquad
 \|h-P(U_{S'})f'\|_2
 <|\T\setminus I|^{1/2}+\eta.
\end{equation}
\end{proposition}

\begin{proof}
We first prepare smooth packets for $f$ and $h$, then put constant tags
on disjoint packets, and finally use Proposition~\ref{prop:prediction}.
Every accuracy appearing below may be made arbitrarily small after the
finite data on which it depends have been fixed.

At a fixed map we use the density estimate
\eqref{eq:density-norm-bound}, which also applies to the nonsmooth
$L^2$ coefficients of an admissible vector. Across a passage, density
preservation instead comes from Lemma~\ref{lem:neutral}.

\paragraph{Converting the angular averages.}
In the current coordinates write
\[
 h=h_\perp+g(u,a),\qquad
 g_0(a)=\int_\T g(u,a)\dd u,\qquad r(u,a)=g(u,a)-g_0(a).
\]
Here $h_\perp$ has zero $\theta$-average, $r$ has zero $u$-average
for each $a$, and $\int g_0=0$.  Fourier truncation and smoothing on the
good arc give a real smooth packet sum $F$ arbitrarily close to $f$.
By \eqref{eq:density-norm-bound}, its density is arbitrarily close to $1$
in $L^1$.

Let $\beta=|\T\setminus I|$.  Fiberwise Parseval gives
\[
 \int_\T|h_\perp(\theta,u,a)|^2\dd\theta
 \le \int_\T|h(\theta,u,a)|^2\dd\theta\le1.
\]
Thus cutting $h_\perp$ off smoothly inside $I$, and then retaining
finitely many of its nonzero angular modes, loses at most
$\sqrt\beta+\eps$ in norm, for arbitrary $\eps>0$.  For example, the
cutoff may be chosen with
$\int_\T(1-\chi)^2\dd u<\beta+\eps^2$.
Denote the resulting packet sum by $H_\perp$.
The simultaneous-list version of Lemma~\ref{lem:neutral}, stated in
Remark~\ref{rem:finite-family-passages}, transports the $F$ and
$H_\perp$ lists through a common passage even when their supports
overlap.  We require density preservation only for the $F$ list; no
orthogonality between the two lists is asserted.

The next passage turns variation in the old $u$ coordinate into
nonzero angular modes; the passage after it does the same for the old
$a$ coordinate. There are two remaining terms, and neither will be
discarded. Before
the first passage integer $N_1$ is chosen, fix a finite nonzero
$u$-Fourier truncation $r_K$ of $r$.  Choose the new good arc in the
$a$ variable and its smooth cutoff $\chi_1(a)$ so that replacing $r_K$
by $\chi_1r_K$ costs arbitrarily little.  These are fixed functions of
the old slow variables.  The cutoff can simultaneously lose arbitrarily
little mass from the finitely many existing packets.  Choose the broad
neutral jitter and then $N_1$, as in Lemma~\ref{lem:neutral}.  In the new
coordinates
\[
 (u,a,b),\qquad b=\theta-H_1(u),
\]
the function $\chi_1(a)r_K(u,a)$ is already a finite sum of nonzero
angular modes supported in the new good arc.  The former packets are
transported with arbitrarily small norm error, and the density of the
$F$ packets changes by arbitrarily little in $L^1$.  Choose $N_1$ also so that every retained transported old index lies
outside $[-K,K]$: its quotient by $N_1$ stays in a fixed compact set
away from zero. This keeps the added modes of $r_K$ disjoint in angular
index from the transported test packets.  The function $g_0(a)$ has been left entirely untouched.

Now fix a finite nonzero $a$-Fourier truncation $(g_0)_L$ of $g_0$.
It has no constant term because $\int g_0=0$.  Apply a second neutral
passage to all existing packets.  Its new coordinates are
\[
 (a,b,d),\qquad d=u-H_2(a).
\]
Choose its good arc and cutoff $\chi_2(b)$ before its final integer
$N_2$.  The term $\chi_2(b)(g_0)_L(a)$ is now a smooth packet sum in
nonzero angular modes.  Its cutoff loses arbitrarily little norm and
cannot introduce an angular constant.  Choose $N_2$ also so that the transported indices lie outside
$[-L,L]$, keeping the newly added modes disjoint from them.  The existing packets are
transported with arbitrarily small norm and density errors.

Relabel the coordinates of the resulting twist as $(\theta,u,a)$.
For every $\eps>0$ the above choices therefore give real smooth packet
sums $F,H$ such that
\begin{equation}\label{eq:prepared-pair}
 \|F-f\|_2<\eps,\qquad
 \|H-h\|_2<\sqrt\beta+\eps,\qquad
 \|p_F-1\|_1<\eps.
\end{equation}
The functions that needed new cutoffs were fixed before the relevant
passage integer was chosen.  In particular, this argument has not
assumed uniform control of coefficients arising only after that choice.

\paragraph{Constant ratios on rectangles.}
Enlarge the finite angular index list so that it contains the modes of
both $F$ and $H$, and write their coefficients as $F_m,H_m$.  For almost
every sufficiently small real $t$,
\begin{equation}\label{eq:ratio-null-set}
 \int_{\{F_m+tH_m=0\}}|H_m|^2\dd u\,\dd a=0
 \quad\hbox{for every }m.
\end{equation}
To see this, at a point where $H_m\ne0$ at most one real $t$ solves
$F_m+tH_m=0$; Fubini applied to $|H_m|^2\dd u\,\dd a\,\dd t$ proves
the assertion.  There are only finitely many indices.  Replace $F$ by
$F+tH$ with such a sufficiently small real $t$.  Its norm and density
changes are arbitrarily small by \eqref{eq:density-norm-bound}, and
reality is preserved.

Define $R_m=H_m/F_m$ where $F_m\ne0$, and zero otherwise.  In view of
\eqref{eq:ratio-null-set}, $H_m=R_mF_m$ almost everywhere.  With
$\dd\nu_m=|F_m|^2\dd u\,\dd a$, one has
$\int|R_m|^2\dd\nu_m=\|H_m\|_2^2$.  Truncating $R_m$ to a bounded
function therefore changes $R_mF_m$ arbitrarily little in $L^2$.
Finite rectangular step functions are dense in $L^2(\nu_m)$: continuous
functions are dense for this finite measure, and uniform continuity
approximates a continuous function by a fine rectangular grid.  Grid
boundaries have $\nu_m$-measure zero, since $\nu_m$ is absolutely
continuous.

Approximate each bounded ratio by such a step function, and replace the
indicators of the finitely many grid rectangles by smooth interior
cutoffs.  Trim enough to obtain positive gaps between distinct retained
rectangles, but lose arbitrarily little in both the $F$ norm and the
tagged norm.  This is possible after the finitely many step values are
known.  Where a ratio has been truncated, a zero step value is allowed;
we do not need to remove the corresponding $F$ mass.  Performing the
construction at positive indices and reflecting it at negative indices
gives packets $F_r$ and constants $b_r$ with conjugate symmetry, such
that
\[
 F^{\mathrm{tag}}=\sum_r F_r\quad\hbox{approximates }f,
 \qquad
 H^{\mathrm{tag}}=\sum_r b_rF_r\quad\hbox{approximates }h
\]
with the errors in \eqref{eq:prepared-pair}, up to arbitrarily small
additions.  Different rectangles in $(m,u,a)$ are disjoint with gaps;
each is its own color.  The density of $F^{\mathrm{tag}}$ remains
arbitrarily close to $1$ in $L^1$.  The number
\[
 M=\max\bigl(1,\max_r|b_r|\bigr)
\]
is finite.  All subsequent approximation errors for tagged vectors are
chosen after $M$ is known.

\paragraph{Separation by angular index alone.}
Two further neutral passages supply exactly the hypotheses of
Proposition~\ref{prop:prediction}.  Here are the support details.
Write the single old angular index of rectangle $r$ as $m_r$, and its
compact coefficient support as $K_r$.  Choose $c$ so that different
$m_r$ have separated ranges, and form
\[
 E_r=\{(m_r(c+\phi(u)),a):(u,a)\in K_r\}.
\]
The compact sets $E_r$ are pairwise disjoint: for a fixed $m_r$ the map
$(u,a)\mapsto(m_r(c+\phi(u)),a)$ is injective, and different indices
were separated by $c$.  Choose smooth cutoffs $\xi_r(s,a)$ supported
in disjoint neighborhoods of these sets and equal to one near them.
On the first passage retain the exact coefficients
\begin{equation}\label{eq:first-separation}
 C_{j,r}(a)=\chi(a)\xi_r(j/N,a)
 \int_\T A_r(u,a)\ee(m_rH(u)-ju)\dd u.
\end{equation}
Here $\chi$ is the new good-arc cutoff; its loss was made small before
choosing the broad neutral jitter and $N$.  Lemma~\ref{lem:passage}
shows that the support restrictions delete arbitrarily little norm as
$N$ increases.  There are finitely many retained $j$, all nonzero.  In
coordinates $(u,a,b)$ the resulting packets are
\[
 C_{j,r}(a)\ee(ju)\ee(m_rb).
\]
Their moduli are independent of the new spare coordinate $b$.
Different old indices cannot share $j$; at a shared $j$, distinct colors
have disjoint compact $a$ supports with positive gaps.

Fix this finite output before making the second passage.  Choose its
$c_2$ to separate the ranges belonging to distinct $j$.  For each fixed
$j$, strict monotonicity of the current speed $\sigma$ on its good arc
makes the compact sets
\[
 Q_{j,r}=\{j(c_2+\sigma(a)):a\in\supp C_{j,r}\}
\]
disjoint as $r$ varies.  Thus all nonempty $Q_{j,r}$ have disjoint open
neighborhoods.  Restrict the second-passage coefficient at new index
$k$ by a smooth cutoff of $k/N_2$ supported in its designated
neighborhood and equal to one near $Q_{j,r}$.  Taking $N_2$ last makes
the deleted norm arbitrarily small, while the retained finite
$k$-index sets are \emph{exactly} disjoint for distinct pairs $(j,r)$.
An individual output term has the form
\[
 B_{k,j,r}(b)\ee(ka)\ee(jd),\qquad d=u-H_2(a),
\]
so its modulus is independent of the final spare coordinate $d$.
No interference can spoil this property, because different pairs have
disjoint final index sets.

Every integral and cutoff used here commutes with the constant tag
$b_r$.  Consequently, if the sum of the individual color transport
errors is less than $\varepsilon/(1+M)$, both the untagged and the
tagged physical errors are less than $\varepsilon$.  Neutral density
preservation and \eqref{eq:density-norm-bound} keep the total density
arbitrarily close to $1$.  At this point the packet colors are disjoint
in angular indices, their moduli are independent of the spare
coordinate, and their mass can be required to be at most $3/2$.

\paragraph{Prediction and whitening.}
Apply Proposition~\ref{prop:prediction} to the tags $b_r/M$, with
prediction tolerance divided by $M$.  Its physical transport errors
are likewise chosen with the factor $M$ included.  Write $G$ for the
final real packet sum, $p$ for its density, and $D(U)G$ for the
resulting approximation to the tagged vector, with the factor $M$
restored.  Here $U$ is the Koopman operator of the final intermediate
map.  We can ensure
\begin{equation}\label{eq:prediction-ready}
 \|G-f\|_2\longrightarrow0,
 \qquad
 \|h-D(U)G\|_2<\sqrt\beta+\eta/2,
 \qquad
 \int_\T\sqrt p\,\dd z\longrightarrow1
\end{equation}
as the preparation and prediction tolerances tend to zero.  For the
last assertion, if the input density is $p_0$, then
\[
 \left|\int\sqrt{p_0}-1\right|
 \le\|\sqrt{p_0}-1\|_2
 \le\|p_0-1\|_1^{1/2},
\]
and Proposition~\ref{prop:prediction} gives an arbitrarily small loss
in $\int\sqrt p$.  Its upper limit is at most one by
$\int\sqrt p\le\|G\|_2$ and $\|G-f\|_2\to0$.

Reality implies $p(-z)=p(z)$ almost everywhere.  Set $E=\{p=0\}$,
choose the flat real vector $v$ of Lemma~\ref{lem:flat-vector} for this
last map, and define
\begin{equation}\label{eq:whitening}
 f'=\left(\frac{\1_{E^c}}{\sqrt p}\right)(U)G+\1_E(U)v.
\end{equation}
The first term is well defined in $L^2$, with squared norm $|E^c|$.
The terms have disjoint spectral supports and densities $\1_{E^c}$
and $\1_E$, respectively.  Indeed, $\1_E(U)G=0$, even if $v$ uses some
of the same angular modes.  The total density of $f'$ is exactly one.
Both multipliers preserve reality, nonzero angular indices, and compact
support in the good arc.  Hence the resulting pair is admissible.
Directly, including the contribution from $E$,
\begin{equation}\label{eq:whitening-distance}
 \|f'-G\|_2^2
 =\int_\T(\sqrt p-1)^2\dd z
 =\|G\|_2^2+1-2\int_\T\sqrt p\,\dd z.
\end{equation}
This tends to zero by \eqref{eq:prediction-ready}, and therefore
$\|f'-f\|_2$ can be made less than the independently specified $\tau$.

Whitening also preserves the prediction.  In the measurable spectral
calculus,
\[
 G=(\sqrt p)(U)f',\qquad
 D(U)G=(D\sqrt p)(U)f'.
\]
Since $f'$ has density one and $D(U)G\in L^2$, the function
$D\sqrt p$ belongs to $L^2(\T,\dd z)$.  Trigonometric polynomials are
dense in this space, so $D(U)G$ belongs to $\calC(U,f')$ and is
approximated by a finite Laurent polynomial in $U$ applied to $f'$.
Choose that approximation with error less than $\eta/2$.  This proves
the last inequality in \eqref{eq:stage-conclusion}; it asserts an
approximation to $h$, not exact membership of $h$ at this finite stage.

Finally, each neutral passage and Proposition~\ref{prop:prediction}
allows an arbitrarily small physical perturbation of both map and
inverse.  There are finitely many such operations here, so they can
all be kept in the prescribed neighborhood.  The last new good arc
may have complement less than $\rho$.  All choices of small physical
errors are made after the finite tags, cutoffs, and coordinate changes
that affect their bounds have been fixed.  Thus neither large tags nor
large derivatives of an earlier coordinate change obstruct the
arbitrary closeness claimed in the proposition.
\end{proof}

\begin{proof}[Proof of Theorem~\ref{thm:main}]
Choose a smooth periodic speed $\phi_0$ which agrees with an increasing
affine function on a nonempty closed interval inside $(0,1)$, and let
$I_0$ be an open subinterval where $\phi_0'>0$.  Start with the twist
$S_0$ and its flat vector $f_0$ from Lemma~\ref{lem:flat-vector}.
Choose a countable family of real smooth centered functions, bounded
by one, whose real linear span is dense in real $L^2_0(\T^3)$; the
physical torus sines and cosines will do.  List them as $(h_n)_{n\ge1}$
so that each occurs infinitely often.

We record the continuity needed to retain approximations.  If
volume-preserving maps $T_i$ and their inverses converge uniformly to
$T$ and $T^{-1}$, then for every fixed $k\in\Z$,
\begin{equation}\label{eq:strong-koopman}
 U_{T_i}^kg\longrightarrow U_T^kg\quad\hbox{in }L^2
 \quad(g\in L^2).
\end{equation}
Uniform convergence of the corresponding fixed iterates follows by
composition on the compact torus.  For continuous $g$ this proves
uniform convergence of the compositions; for general $g$, approximation
by continuous functions and unitarity prove \eqref{eq:strong-koopman}.
Consequently, for every fixed Laurent polynomial $P$ the map
\[
 (T,g)\longmapsto P(U_T)g
\]
is continuous into physical $L^2$ when maps and inverses have the
$C^\infty$ topology and vectors have the $L^2$ topology.  Explicitly,
if $P(w)=\sum_k c_kw^k$, its difference is bounded by
\[
 \sum_k|c_k|\|g_i-g\|_2+
 \sum_k|c_k|\|U_{T_i}^kg-U_T^kg\|_2.
\]

Inductively apply Proposition~\ref{prop:stage} with $h=h_n$ and
$\eta=2^{-n}$, obtaining $(S_n,f_n)$, a good arc $I_n$, and a finite
Laurent polynomial $P_n$, such that
\begin{equation}\label{eq:stage-polynomial}
 |\T\setminus I_n|<4^{-n},\qquad
 \|h_n-P_n(U_{S_n})f_n\|_2
 <|\T\setminus I_{n-1}|^{1/2}+2^{-n}.
\end{equation}
Use the arbitrary closeness in that proposition also to impose
\[
 \|f_n-f_{n-1}\|_2<2^{-n},
 \qquad
 d_k(S_n,S_{n-1})+d_k(S_n^{-1},S_{n-1}^{-1})<2^{-n}
 \quad(0\le k\le n),
\]
where $d_k$ are ambient $C^k$ sup-norm distances induced by a fixed
smooth embedding of the torus. At the same time impose
the finitely many open conditions
\begin{equation}\label{eq:retain-polynomial}
 \|P_j(U_{S_n})f_n-P_j(U_{S_{n-1}})f_{n-1}\|_2<2^{-n}
 \qquad(1\le j<n).
\end{equation}
Their compatibility follows from the continuity just proved.  No bound
uniform in $j$ on the degrees or coefficients of $P_j$ is required.

For each fixed $k$, the map and inverse sequences have summable tails
in $C^k$.  They therefore converge smoothly to maps $T$ and $R$.
Uniform convergence, compactness, and
$S_nS_n^{-1}=S_n^{-1}S_n=\operatorname{id}$ give
$TR=RT=\operatorname{id}$, so $T$ is a smooth diffeomorphism with
smooth inverse.  For every continuous $g$,
\[
 \int g\circ T=\lim_n\int g\circ S_n=\int g,
\]
so normalized volume is invariant.  The vectors $f_n$ converge in
$L^2$ to a real centered vector $f$.

Every completed stage has density exactly one.  Combining
\eqref{eq:strong-koopman} with $f_n\to f$, for all $r,s\in\Z$ we obtain
\begin{equation}\label{eq:limit-orthonormal}
 \langle U_T^rf,U_T^sf\rangle
 =\lim_n\langle U_{S_n}^rf_n,U_{S_n}^sf_n\rangle
 =\delta_{rs}.
\end{equation}
The retained finite polynomials give, for every $j$,
\[
 \|h_j-P_j(U_T)f\|_2
 \le |\T\setminus I_{j-1}|^{1/2}+2^{-j}
       +\sum_{n>j}2^{-n}
 =|\T\setminus I_{j-1}|^{1/2}+2^{1-j}.
\]
Along the infinitely many occurrences of any fixed test the right side
tends to zero.  Hence every test belongs to $\calC(U_T,f)$.  Their
complex linear span is dense in complex $L^2_0(\T^3)$, whereas every
orbit vector is centered.  Thus
\[
 \calC(U_T,f)=L^2_0(\T^3).
\]
Together with \eqref{eq:limit-orthonormal}, this says that the bilateral
orbit of $f$ is an orthonormal basis of the entire centered space.
With $w=\ee(z)$, the assignment $U_T^nf\mapsto w^n$ extends to a
unitary map onto $L^2(S^1,m)$ intertwining $U_T$ with multiplication
by $w$.

Multiplication by $w$ has no nonzero invariant vector, since
$(w-1)g(w)=0$ forces $g=0$ almost everywhere.  Consequently every
$T$-invariant $L^2$ function is constant, which proves ergodicity.
The manifold $\T^3$ is compact, connected, smooth, and without
boundary, and its normalized flat volume has strictly positive smooth
density.  These are all the properties required in
Theorem~\ref{thm:main}.
\end{proof}

\section{Dynamical consequences}
\label{sec:consequences}

The spectral model gives ordinary mixing directly. The passage to higher
orders uses the separate multiple-mixing theorem; the entropy and Lyapunov
conclusions do not depend on that additional input.

\begin{corollary}\label{cor:dynamical-profile}
For the diffeomorphism $T$ and normalized standard volume $\mu$ of
Theorem~\ref{thm:main}, every integer $k\ge2$ and every fixed choice of
measurable sets $A_1,\ldots,A_k\subseteq\T^3$ satisfy
\[
 \mu\left(\bigcap_{i=1}^k T^{-t_i}A_i\right)
 \longrightarrow\prod_{i=1}^k\mu(A_i)
\]
whenever $t_1<\cdots<t_k$ are integers and
$\min_{1\le i<k}(t_{i+1}-t_i)\to\infty$. Moreover, the
Kolmogorov--Sinai entropy satisfies $h_\mu(T)=0$, and all three Lyapunov
exponents of $T$, counted with multiplicity, vanish at $\mu$-almost every
point.
\end{corollary}

\begin{proof}
Let $V:L^2_0(\T^3,\mu)\to L^2(S^1,m)$ be the unitary spectral model in
Theorem~\ref{thm:main}. For $a,b\in L^2_0(\T^3,\mu)$,
\[
 \int_{\T^3}(a\circ T^n)\overline b\dd\mu
 =\int_{S^1}w^n(Va)(w)\overline{(Vb)(w)}\dd m(w)
 \longrightarrow0\qquad(|n|\to\infty).
\]
Indeed, $(Va)\overline{Vb}\in L^1(m)$ by Cauchy--Schwarz, so the limit is
the Riemann--Lebesgue lemma. Taking centered indicators gives ordinary
mixing. The map $T$ is an invertible probability-preserving transformation
with measurable inverse, so the multiple-mixing theorem
\cite[Theorem~1.1]{MultipleMixing2026} gives every order $k\ge3$.
Invariance shifts $t_1$ to zero and identifies its positive time gaps with
$t_{i+1}-t_i$; the case $k=2$ is ordinary mixing.

The full Koopman operator is the direct sum of the one-dimensional
constants and the simple Lebesgue model on all of $L^2_0$, and hence has
finite spectral multiplicity. Normalized torus volume is a Lebesgue
probability space, and $T$ is a measure-preserving automorphism.
Rokhlin's finite-multiplicity entropy theorem
\cite[\S 7, p.~13]{Rokhlin1960} therefore gives $h_\mu(T)=0$.

The derivative and inverse derivative are bounded on the compact torus,
so the invertible Oseledets integrability conditions hold. Write
$\lambda_1(x),\lambda_2(x),\lambda_3(x)$ for the exponents repeated
according to multiplicity at regular points. Pesin's entropy formula,
as stated and reproved by Ma\~n\'e \cite[p.~95, equation~(2)]{Mane1981},
applies to the $C^\infty$ diffeomorphism and its absolutely continuous
invariant probability $\mu$. Thus
\[
 0=h_\mu(T)=\int_{\T^3}\sum_{i=1}^3\max\{\lambda_i(x),0\}\dd\mu(x).
\]
The nonnegative integrand vanishes almost everywhere, so every exponent
is nonpositive there. Preservation of standard volume gives
$|\det DT_x|=1$, and at every Oseledets-regular point
\[
 \sum_{i=1}^3\lambda_i(x)
 =\lim_{n\to\infty}\frac1n\log|\det D(T^n)_x|=0.
\]
The three nonpositive exponents must therefore all be zero at
$\mu$-almost every point.
\end{proof}

\begingroup
\small
\bibliographystyle{plain}
\bibliography{references}
\endgroup
\end{document}
